\documentclass[IJPHM, 2026, 0]{PHMSociety}

% Declare Packages
\usepackage[T1]{fontenc}
\usepackage{graphicx}
\usepackage{amsmath}
\usepackage{amssymb}
\usepackage{algorithm}
\usepackage{algpseudocode}
\usepackage{booktabs}
\usepackage{multirow}
\usepackage{array}
\usepackage{url}
\urlstyle{tt}
% Typeset a long URL one size smaller so it fits on a single line inside the
% narrow abstract column instead of breaking mid-word at a hyphen.
\newcommand{\nburl}[1]{{\footnotesize\url{#1}}}
\usepackage{placeins}
\usepackage{float}

% Float placement: allow denser float pages and stop LaTeX from deferring
% tables/figures into a cluster at the end of the document.
\renewcommand{\topfraction}{0.9}
\renewcommand{\bottomfraction}{0.8}
\renewcommand{\textfraction}{0.07}
\renewcommand{\floatpagefraction}{0.7}
\setcounter{topnumber}{4}
\setcounter{bottomnumber}{4}
\setcounter{totalnumber}{8}

% Single-author license footnote: override the class default, which hard-codes
% "et al." after the name (PHMSociety.cls). Text otherwise verbatim.
\renewcommand{\phmLicenseFootnote}[1]{%
\insert\footins{\noindent\footnotesize{#1. This is an open-access article distributed under the terms of the Creative Commons Attribution 3.0 United States License, which permits unrestricted use, distribution, and reproduction in any medium, provided the original author and source are credited.}}}

\begin{document}

% Paper Title
\title{A Latch-On-Resistant, False-Alarm-Gated Lead-Time Metric for Bearing Anomaly Detection --- and What It Reveals About SCADA-Rate Logging}

% Authors List
\author{%
	Ansh Goyal\authorNumber{1}
}

% Author Affiliations
\address{%
	\affiliation{1}{Jaypee Institute of Information Technology, Noida, India}{%
		{\email{anshgoyal500@gmail.com}}
		}
}

% Create the title
\maketitle
\pagestyle{fancy}
\thispagestyle{plain}

% PHM Society Distribution License Information
\phmLicenseFootnote{Ansh Goyal}

% Abstract
\begin{abstract}%
Anomaly detectors for rolling-element-bearing prognostics are typically evaluated by classification accuracy, which answers whether a fault is detected but not how early or at what false-alarm cost. We argue that lead time---the warning interval before failure---is the operationally decisive quantity, and we make it the primary evaluation object through a latch-on-resistant, false-alarm-gated metric. Warning time is credited only when the pre-onset false-alarm rate lies within a stated budget $\tau$, so a degenerate constant-on detector is provably scored invalid. We then ask whether the bin-averaged logging of a SCADA historian destroys detection lead time relative to simple decimation. Eleven detectors---four SPC charts, Isolation Forest, three deep reconstruction autoencoders, an RMS-trend baseline, and two one-class models---are assessed under a controlled sampling sweep across four inferential run-to-failure datasets: XJTU-SY ($n=10$), FEMTO/PRONOSTIA ($n=6$), University of Ferrara ($n=6$), and NASA IMS ($n=3$). On IMS the variance-sensitive charts show a positive median shift ($3\sigma$ $+15.1$~h), but the direction is not consistent across the three runs and does not survive Holm correction across the $N=44$ family (smallest adjusted $p=1.00$). We state the cross-dataset conclusion as a bounded equivalence rather than an accepted null: against a pre-specified $\pm 1$~h operational margin, the run-level aggregate-minus-decimate difference is equivalent to zero in all 30 dataset $\times$ detector cells on the three multi-bearing campaigns, while on IMS the intervals are too wide at $n=3$ to decide in 8 of 10 cells and favour aggregation beyond the margin in the other 2, and the single-asset gas-turbine record (ONGC, $n=1$) cannot be tested. No cell shows decimation superior beyond the margin; the one place decimation leads, Isolation Forest on FEMTO ($-0.26$~h, 95\% CI $[-0.54,-0.06]$), lies inside it. The full pipeline, detectors, and result tables are released at \nburl{https://github.com/Ansh-Goyal01/Scada-Leadtime-Benchmark}.
\end{abstract}

\medskip
\noindent\textbf{Keywords:} bearing anomaly detection, lead-time metric, false-alarm gating, SCADA historian, prognostic horizon, statistical process control, conformal prediction
\medskip


\section{Introduction}
\label{sec:intro}
Condition-based maintenance (CBM) and prognostics and health management (PHM) promise to replace fixed-interval servicing with data-driven intervention, reducing both unplanned downtime and unnecessary maintenance \cite{jardine2006review,lee2014prognostics,lei2018machinery}. Rolling-element bearings are the most common failure mode in rotating machinery, and a large literature applies statistical control charts, classical machine learning, and deep models to detect incipient bearing faults from vibration \cite{randall2011rolling,zhao2019deep}. The economic case is well established: a warning delivered hours or days before a catastrophic seizure lets an operator schedule an orderly shutdown, order parts, and avoid secondary damage to shafts, housings, and couplings.

Despite this, most published detectors are evaluated with classification metrics---accuracy, precision/recall, ROC--AUC. These answer ``was the failure detected?'' but not the question a maintenance planner actually asks: ``how many hours of warning did we get, and how often did the system cry wolf?'' A detector that warns 48~h ahead at a low nuisance-alarm rate is operationally distinct from one that fires minutes before breakdown or one that alarms every day. The prognostics community has long advocated time-aware metrics \cite{saxena2008metrics,saxena2010metrics,goebel2011comparison}, yet lead time is rarely the headline quantity in detector benchmarks, and ad hoc lead-time definitions are easy to game. In particular, any lead-time metric that does not pair warning time with a false-alarm constraint can be trivially maximized by a detector that simply alarms from the first sample---a ``latch-on'' exploit that earns the full run length as lead time while being operationally useless.

A second, eminently practical question motivates this study. Industrial SCADA historians rarely store raw high-rate vibration; they store decimated or time-averaged summaries at a coarse logging interval \cite{tautzweinert2017using}---seconds to minutes, rather than the tens of kilohertz at which accelerometers are sampled. The folklore concern is that averaging destroys the kurtosis and crest-factor transients that signal an incipient spall: if true, the very act of long-term storage would blind any prognostic algorithm that consumes the historian. Whether this is actually true is an empirical question that, to our knowledge, has not been tested with a sound lead-time metric across multiple datasets. The answer matters: if averaging is harmless, operators can store compact, cheap, bin-averaged streams without sacrificing warning time; if it is harmful, the implication is that prognostics-grade monitoring requires separate high-rate acquisition.

This paper makes both questions precise and answers them. We state the arc explicitly:

\textbf{Problem.} Detectors are ranked by classification accuracy, which ignores how early a warning arrives and at what false-alarm cost; and it is unknown whether the coarse, bin-averaged logging of a SCADA historian destroys that warning time relative to plain decimation.

\textbf{Solution.} We make lead time gated by a leakage-free false-alarm rate the primary metric---warning time counts only when the pre-onset false-alarm rate is within budget---so that the degenerate latch-on detector is provably scored invalid rather than winning.

\textbf{Method.} Under this metric we run a controlled sampling sweep that holds window content and alarm persistence constant, applied to nine detectors (four SPC charts, Isolation Forest, three deep reconstruction autoencoders, and an RMS-trend baseline; Deep SVDD additionally evaluated) across four inferential run-to-failure datasets---XJTU-SY ($n=10$), FEMTO/PRONOSTIA ($n=6$), University of Ferrara ($n=6$), and NASA IMS ($n=3$), plus a real single-asset ONGC gas-turbine SCADA record analyzed as a case study (Appendix~D.2)---with run-level significance tests (the run, not the within-run sampling factor, is the unit of inference) and Holm correction.

\textbf{Finding.} The folklore is refuted on every dataset tested: historian-style aggregation does not destroy lead time. On the large multi-bearing sets (XJTU-SY, FEMTO, Ferrara) the difference is negligible; on IMS it shows a consistent positive trend (not significant at $n=3$); the single-asset turbine record agrees (case study, Appendix~D.2). The defensible cross-dataset conclusion is non-destruction---decimation is never better, and averaging is at worst neutral.

\textbf{Contributions.} We make five contributions.

\begin{enumerate}
\item \emph{A latch-on--resistant, false-alarm-gated, onset-relative lead-time metric.} Warning time is credited only when the false-alarm rate over a data-anchored, leakage-free pre-onset region is within a stated budget. We prove---by construction and by a unit test that is part of the released code---that a constant-on (latch-on) detector is scored invalid, while a detector that fires just after the true degradation onset scores valid at zero false-alarm rate. We do not claim the metric is unconditionally ungameable; we discuss residual gaming surfaces in Section~\ref{sec:limitations}.
\item \emph{A formal proof that gating changes the answer.} We give, to our knowledge, the first formal comparison showing that an ungated prognostic-horizon metric (the Saxena PH) produces a different and operationally misleading detector recommendation relative to the false-alarm-budgeted metric (Section~\ref{sec:saxena}, Tables~\ref{tab:latchon} and \ref{tab:phrank}): on IMS the detectors PH ranks first by raw lead---Hotelling $T^2$, Isolation Forest, and EWMA---each have no valid operating point at a 10\% false-alarm budget, whereas our metric correctly identifies $3\sigma$ as the deployable choice.
\item \emph{A controlled SCADA-sampling methodology} that holds window content and alarm persistence constant, so that any lead-time difference between the historian's aggregate mechanism and simple decimation reflects genuine information loss rather than a window-counting artifact of coarser time grids.
\item \emph{A multi-dataset null result that is statistically honest by construction.} Across XJTU-SY ($n=10$), FEMTO/PRONOSTIA ($n=6$), University of Ferrara ($n=6$), and NASA IMS ($n=3$), with bootstrap confidence intervals and run-level significance tests that treat the run---not the within-run sampling factor---as the unit of inference, SCADA-rate aggregation does not destroy lead time---and on IMS shows a consistent positive trend for the variance-sensitive control charts---while decimation is never better; nothing survives Holm correction across the $N=40$ family.
\item \emph{Validation:} wiring and empirical validation of a conformal detector for distribution-free false-alarm control, a lead-time--vs--false-alarm trade-off analysis, feature-group and spectral ablations that localize the prognostic signal to time-domain statistics on IMS, and a characterization of deep reconstruction models (LSTM-AE, TCN-AE, Transformer-AD) showing they do not become competitive with SPC charts within the tested training-fraction range $[0.20, 0.60]$ on short-lived bearings; whether a crossover exists on long-runway assets remains open (Section~\ref{sec:mintrain}).
\end{enumerate}

\textbf{Scope of claims.} The empirical results are for rolling-element bearings under approximately steady operating conditions; we make no claim of transfer to gearboxes, blades, or variable-speed machinery. The sampling question is specifically bin-averaged historian storage versus decimated storage. The metric itself is failure-mode-agnostic: it applies wherever a run-to-failure trajectory and a stationary normal region exist.

We deliberately report a result that is not the one we set out to find. Our working hypothesis was that averaging would smear incipient-fault transients and shorten lead time. It does not. We regard reporting the refutation, with its dataset-dependent direction made explicit, as more valuable than a tidy but unsupported positive claim.


\section{Related Work}

\subsection{Prognostics, CBM, and Lead-Time Metrics}
Reviews of machinery diagnostics and prognostics \cite{jardine2006review,lee2014prognostics,lei2018machinery} and of remaining-useful-life (RUL) estimation \cite{si2011remaining} establish the CBM/PHM framing in which a health indicator is tracked and an alarm or RUL estimate is produced. \citeA{saxena2008metrics}, \citeA{saxena2010metrics}, and \citeA{goebel2011comparison} propose a family of prognostic-performance metrics---prognostic horizon, $\alpha$--$\lambda$ accuracy, relative accuracy, and timeliness---that explicitly reward early, well-calibrated predictions. Our onset-relative lead time is in this spirit: it measures warning time relative to a degradation event rather than to an arbitrary positional marker. We depart from prior practice in one key respect: we make the false-alarm budget an integral part of the metric's validity condition rather than a separately reported number, which is precisely what closes the latch-on loophole. We further provide, to our knowledge, the first formal comparison showing that an ungated prognostic-horizon metric produces a different and operationally misleading detector ranking relative to a false-alarm-budgeted metric (Section~\ref{sec:saxena}, Tables~\ref{tab:latchon} and \ref{tab:phrank}). The disagreement is not marginal: on IMS the detectors the prognostic horizon ranks first by raw lead---Hotelling $T^2$, Isolation Forest, and EWMA---each have no valid operating point at a 10\% false-alarm budget, while our metric correctly identifies $3\sigma$ as the deployable choice.

\subsection{Anomaly Detectors and Control Charts}
Univariate and multivariate statistical-process-control charts---the $3\sigma$ rule, the exponentially weighted moving average (EWMA) \cite{roberts1959control}, and Hotelling's $T^2$ \cite{hotelling1947multivariate}---remain workhorses of industrial monitoring \cite{montgomery2009introduction} because they are interpretable and require little training data. In the unsupervised machine-learning family, Isolation Forest \cite{liu2008isolation} and the one-class SVM \cite{scholkopf2001estimating} are widely deployed; deep models such as LSTM autoencoders \cite{hochreiter1997long,malhotra2016lstm,hundman2018detecting} and Deep SVDD \cite{ruff2018deep} have been applied to machine-health monitoring \cite{zhao2019deep}. We evaluate representatives of both the classical and learning families under a single, fixed hyperparameter protocol so that the comparison is about the evaluation and the sampling question, not about detector tuning.

\emph{Relation to published FEMTO/XJTU results.} The FEMTO/PRONOSTIA and XJTU-SY campaigns are best known through the IEEE PHM 2012 challenge and subsequent remaining-useful-life (RUL) studies \cite{wang2020hybrid,nectoux2012pronostia}, whose leading methods report a RUL-error score on the deliberately truncated Test/Validation copies of the bearings. Those numbers are not directly comparable to ours for two structural reasons: (i) they estimate a continuous RUL on runs with no observed failure, whereas our metric measures early-warning lead time on the complete run-to-failure trajectory and gates it on a pre-onset false-alarm budget; and (ii) they do not report a false-alarm rate at all, so a method's RUL score says nothing about whether its alarm would survive our $\tau$ gate---the very quantity we argue is decisive. We therefore do not tabulate their RUL scores against our lead times, which would be a category error. The classic normalized-RMS health-indicator with a first-prediction-time threshold---the workhorse baseline underlying most of that literature---is represented here by our rms\_trend detector, evaluated under the same gated metric as every other method; re-evaluating specific published architectures under the gated metric (which requires their original code) is left as future work.

\subsection{Uncertainty Quantification and Conformal Prediction}
Conformal prediction \cite{vovk2005algorithmic,laxhammar2014online} provides a distribution-free guarantee that the false-alarm rate is bounded by a target $\alpha$ under exchangeability. This is attractive for prognostics because it converts a detector's raw score into an alarm with an interpretable, finite-sample false-alarm guarantee. We wire conformal calibration into the lead-time pipeline and, rather than assuming the guarantee, test it empirically across runs---including a run where exchangeability provably fails because the ``normal'' region already contains degradation.

\subsection{Vibration Features}
Bearing-defect frequencies, spectral kurtosis \cite{antoni2006spectral}, and envelope (demodulation) analysis \cite{randall2011rolling} are the classical frequency-domain tools for bearing diagnostics. We include FFT band energies, spectral kurtosis, and Hilbert-envelope amplitudes at the bearing characteristic frequencies, and we ablate them against pure time-domain statistics to see which family actually carries prognostic lead time under our gated metric.

\subsection{SCADA-Based Condition Monitoring}
\citeA{tautzweinert2017using} review SCADA-data condition monitoring and note explicitly that historian aggregation is a recognized but under-quantified concern: practitioners suspect that coarse averaging hurts fault detection, but the magnitude of the effect has not been measured with a principled warning-time metric. The present paper addresses exactly this gap, and reaches the perhaps surprising conclusion that, for bearing faults, the concern is not borne out.


\section{Datasets}
\label{sec:datasets}
We use four inferential run-to-failure datasets spanning controlled-lab and cross-condition regimes, together with a single-asset real-historian case study (Table~\ref{tab:datasets}). They were chosen so that the central sampling claim is tested under genuinely different conditions: a high-rate raw-waveform rig where aggregation can be simulated faithfully (IMS), two multi-bearing campaigns (XJTU-SY and FEMTO/PRONOSTIA) that supply the statistical power a three-run rig cannot, a second modern constant-condition rig (Ferrara), and---as a case study only---a real historian stream that already is SCADA-rate data (ONGC, Appendix~D.2).

\emph{IMS.} The NASA/IMS bearing dataset \cite{qiu2006wavelet} provides raw 20.48~kHz accelerometer waveforms for three run-to-failure tests on a shaft loaded with four bearings. Because the waveforms are raw, the aggregate/decimate sweep can be applied faithfully at the signal level, and frequency-domain features (spectral kurtosis, envelope) are meaningful. IMS test 3's failure label is set as documented in the released repository; the justification appears in Appendix~C.

\emph{XJTU-SY.} The XJTU-SY dataset \cite{wang2020hybrid} contributes ten run-to-failure bearings under two operating conditions (five bearings each) at 25.6~kHz, sampled in 1.28~s snapshots once per minute. With $n=10$ it provides the strongest generalization test in this study; its published lifetimes range widely (52--533~min), which deliberately stresses the metric with both long-degrading and abrupt-failure bearings.

\emph{FEMTO/PRONOSTIA.} The FEMTO/PRONOSTIA platform \cite{nectoux2012pronostia} contributes an independent run-to-failure campaign: six Learning\_set bearings under three operating conditions, with horizontal and vertical accelerometers sampled in 0.1~s snapshots at 25.6~kHz once every 10~s. We use only the six Learning\_set bearings because they are run until the platform's 20~g vibration stop---a failure, and hence a lead time, is actually observed---whereas the Test\_set/Validation\_set copies are deliberately truncated mid-life for the IEEE PHM 2012 RUL-prediction challenge and would make lead time undefined. Lifetimes span ${\sim}1.4$--7.8~h. We deliberately exclude the Paderborn dataset \cite{lessmeier2016condition} from the lead-time study: it comprises pre-damaged bearings measured at fixed operating conditions for fault classification, with no degradation timeline and no failure event, so it cannot test a lead-time or sampling claim. Using it for the headline would be a category error, and we say so rather than pad the dataset count.

\emph{University of Ferrara.} The University of Ferrara dataset \cite{arpa2024ferrara} contributes a fully independent run-to-failure campaign: six self-aligning double-row ball bearings under constant speed and radial load, with continuous uniaxial vibration recorded at 25.6~kHz throughout each bearing's entire lifecycle (five-second snapshots, contiguous). Lifetimes span 0.7--6.8~h, similar to FEMTO. Operating conditions are fixed, making the dataset directly compatible with our onset detector without normalization; we use all six available runs (E1--E6). The single accelerometer channel is duplicated into the two channel slots so the channel-invariant feature schema applies unchanged. The dataset is included to further test the non-destruction hypothesis with an independent 2024 rig; like FEMTO and XJTU-SY, its short bearing lifetimes mean this campaign cannot confirm or deny the IMS positive trend.

XJTU-SY ($n=10$) is our primary inferential dataset; IMS ($n=3$) is retained for its raw-waveform access, which permits the faithful aggregate/decimate simulation and the mechanism analysis (Appendix~D.1).

\begin{table}[t] \small
	\begin{center}
	\caption{Datasets used in this study. ``Role'' indicates the function each dataset plays in the argument. The four inferential datasets supply the run-level statistics; ONGC is a single-asset case study (Appendix~D.2) and Paderborn is excluded.}
	\label{tab:datasets}
	\begin{tabular}{l l l c}
		\hline \hline
		\textbf{Dataset} & \textbf{Role} & \textbf{Rate} & \textbf{Runs} \\
		\hline
		IMS & controlled lab & 20.48 kHz raw & 3 \\
		ONGC & case study (App.~D.2) & 10 s historian & 1 \\
		XJTU-SY & generalization & 25.6 kHz, 1/min & 10 \\
		FEMTO & generalization & 25.6 kHz, 1/10 s & 6 \\
		Ferrara & generalization & 25.6 kHz, cont. & 6 \\
		Paderborn & excluded (no RUL) & 64 kHz & -- \\
		\hline \hline
	\end{tabular}
	\end{center}
\end{table}


\section{Methodology}

\subsection{Problem Formulation and Notation}
For each run we observe a sequence of snapshots indexed by time $t$. A feature extractor maps each snapshot (and a short trailing window) to a vector $x_t \in \mathbb{R}^d$. A detector assigns a score $s_t$ and, via a threshold and a persistence filter, an alarm indicator $a_t \in \{0,1\}$. Each run has a failure time $t_f$ and is split at a train boundary $t_b$; all scaling and threshold statistics are fit on $\{t \le t_b\}$ only. The two quantities we care about are the lead time (how long before $t_f$ the first sustained alarm fires) and the pre-onset false-alarm rate (how often the detector alarms while the machine is still healthy). The remainder of this section makes each ingredient precise.

\subsection{Channel-Invariant Feature Schema}
Each snapshot yields per-channel time-domain statistics---RMS, kurtosis, skewness, peak-to-peak, crest factor, and variance---and, for raw-waveform datasets, frequency-domain features: FFT band energies, spectral kurtosis \cite{antoni2006spectral}, and Hilbert-envelope amplitudes at the bearing characteristic frequencies (BPFO, BPFI, BSF, FTF), computed from the dataset geometry.

A naive feature scheme that concatenates per-channel statistics and all cross-channel pairs is channel-count dependent and ill-posed: it produces 445 features for a 4-channel run and 1465 for an 8-channel run, against only ${\sim}100$ training snapshots, i.e.\ $p \gg n$. Such a space overfits, is not comparable across runs with different channel counts, and silently inflates the apparent dimensionality. We therefore aggregate each statistic family across physical channels (taking the mean and the max over channels) and summarize each aggregated series over a rolling window with four descriptors (mean, standard deviation, max, and linear-trend slope), plus one cross-channel RMS-coherence feature. This yields a fixed 49-dimensional space that is identical across 4- and 8-channel assets and directly comparable across datasets. A train-fit variance top-$k$ selection ($k \le 50$) then guarantees $p < n$. All scaling is fit on the training split only, so no test-region statistic leaks into the features.

\subsection{Leakage-Free Degradation Onset}
Lead-time metrics that reference an arbitrary positional marker (e.g.\ ``the first 5\% of the test window is normal'') are gameable, because the marker is not tied to the physics of the run. We instead compute a single degradation onset $t_o$ per run from a scalar health indicator $h(t)$. The default indicator is the element-wise maximum of two baseline-standardized trends---the mean-RMS trend and the mean-kurtosis trend---so that it responds to whichever degradation mode (amplitude growth or impulsiveness) appears first:
\begin{equation}
h(t) = \max\!\big(z_b[\text{RMS}](t),\, z_b[\text{kurt}](t)\big),
\label{eq:health}
\end{equation}
where the baseline mean $\mu_b$ and standard deviation $\sigma_b$ used in the z-score $z_b[\cdot]$ are estimated only on data up to the train boundary $t_b$, so $t_o$ cannot see the test region. With band $\beta = \mu_b + k\sigma_b$ (default $k=4$), the onset is the start of the last sustained excursion above $\beta$ before failure,
\begin{equation}
t_o = \min\{\, t \in R_{\text{last}} : h(t) > \beta \,\},
\label{eq:onset}
\end{equation}
where $R_{\text{last}}$ is the final run of consecutive above-band samples, merging gaps shorter than a tolerance and requiring a minimum persistence. Using the last sustained excursion makes the onset robust to early run-in transients and to the single trailing machine-off sample, both of which can otherwise spuriously mark an early onset. Algorithm~\ref{alg:onset} states the procedure.

\begin{algorithm}
\caption{Leakage-free terminal degradation onset}
\label{alg:onset}
\begin{algorithmic}[1]
\State \textbf{Input:} health series $h(t)$; train boundary $t_b$; failure $t_f$; threshold multiplier $k$; gap/persistence tolerances
\State $\mu_b, \sigma_b \gets$ mean, std of $h$ on $\{t \le t_b\}$
\State $\beta \gets \mu_b + k\,\sigma_b$
\State $E \gets \{t < t_f : h(t) > \beta\}$ \Comment{above-band samples}
\State merge gaps in $E$ shorter than tolerance; drop runs shorter than persistence
\State $R_{\text{last}} \gets$ final surviving run in $E$
\State \textbf{return} $t_o \gets \min R_{\text{last}}$
\end{algorithmic}
\end{algorithm}

The persistence (5 samples) and gap (10 samples) tolerances of Algorithm~\ref{alg:onset} are fixed in samples of the health-indicator series, not physical time, and are held constant across datasets. Because the datasets are logged at different native intervals---${\sim}10$~min per snapshot for IMS (after resampling to a uniform grid), 1~min for XJTU-SY, and 10~s for FEMTO and ONGC---the same 5/10-sample bars correspond to different physical durations: ${\sim}50/100$~min on IMS, 5/10~min on XJTU, and 50/100~s on FEMTO/ONGC. This keeps the onset rule algorithmically identical everywhere at the price of exact physical-time comparability; the difference is negligible for long records but non-trivial for the shortest bearings---XJTU Bearing2\_4 (42~min total life) must clear a ${\sim}5$-minute persistence bar, roughly 12\% of its lifetime---so onset positions on the very shortest bearings carry this caveat.

\medskip
\noindent\textit{Caveat: onset positions on bearings with total lifetime below 2~h are treated as fragile throughout; per-bearing results in Tables~\ref{tab:xjtu}, \ref{tab:femto}, \ref{tab:ferrara}, and \ref{tab:perbearing} flag these cases explicitly rather than silently discarding them.}
\medskip

\noindent The sensitivity analyses probe the threshold multiplier $k$ (Table~\ref{tab:onsetsens}), not these fixed timing tolerances, which we leave in sample units for algorithmic reproducibility.

\subsection{Onset-Relative Lead-Time Metric}
Let the first sustained alarm time be FAT and the failure time $t_f$. The lead time is
\begin{equation}
L = \max\!\left(0,\, \frac{t_f - \text{FAT}}{3600}\right) \text{ hours},
\label{eq:lead}
\end{equation}
with $L=0$ if there is no alarm or if the alarm is at or after failure. The pre-onset false-alarm rate is the fraction of windows strictly before the onset that are flagged,
\begin{equation}
\text{FAR}_{\text{pre}} = \frac{|\{t < t_o : a(t)=1\}|}{|\{t < t_o\}|},
\label{eq:farpre}
\end{equation}
measured over the entire apparently-normal region rather than a positional sample. An alarm is valid if and only if
\begin{equation}
L > 0 \;\wedge\; \text{FAR}_{\text{pre}} \le \tau,
\label{eq:valid}
\end{equation}
with budget $\tau=0.10$ throughout.

\emph{Why the metric defeats the latch-on exploit.} Reporting $L$ and $\text{FAR}_{\text{pre}}$ jointly, with the conjunction in Eq.~(\ref{eq:valid}), defeats the latch-on exploit. A detector that fires from the first sample earns a large $L$ (the full run length) but floods the entire pre-onset region, giving $\text{FAR}_{\text{pre}} \approx 1 > \tau$ and hence invalid. Conversely, a detector that first fires just after $t_o$ has $\text{FAR}_{\text{pre}}=0$ and a positive $L$, scoring valid. A unit test in the released code asserts exactly these two cases: a constant-on detector scores invalid, and an oracle that fires at $t_o^+$ scores valid at zero pre-onset FAR. Because $t_o$ is computed only from training-region statistics (Eq.~\ref{eq:onset}), neither $L$ nor $\text{FAR}_{\text{pre}}$ can be tuned using test-region information. We stress that this defeats the specific latch-on exploit; it is not a claim of unconditional ungameability (see Section~\ref{sec:limitations}).

\subsection{Detectors}
The headline set comprises nine detectors, evaluated under one fixed hyperparameter protocol:
\begin{itemize}
\item \emph{RMS-trend} (naive baseline): alarm when the mean-RMS exceeds $k\sigma$ above its training baseline. Included to verify that the metric correctly penalizes a weak detector.
\item \emph{$3\sigma$ rule:} alarm when the $\ell_2$-norm health index exceeds three training standard deviations.
\item \emph{EWMA} \cite{roberts1959control}: an exponentially weighted moving-average control chart ($\lambda=0.2$, $k=3.0$), sensitive to small sustained shifts.
\item \emph{CUSUM} \cite{page1954continuous}: a two-sided tabular cumulative-sum control chart ($k=0.5$, $h=5.0$) on the monitored magnitude---the classic change-point chart, most sensitive to a small sustained mean shift such as slow degradation onset. Deterministic.
\item \emph{Hotelling $T^2$} \cite{hotelling1947multivariate}: computed in a PCA-reduced subspace (rank $\ll n$) so the covariance is well-conditioned.
\item \emph{Isolation Forest} \cite{liu2008isolation}: an unsupervised ensemble that isolates anomalies with few random splits.
\item \emph{Deep reconstruction baselines:} an LSTM autoencoder \cite{malhotra2016lstm}, a Temporal Convolutional Network autoencoder (dilated causal convolutions), and a lightweight reconstruction Transformer \cite{vaswani2017attention}. Each is trained on normal feature-window sequences only and scores a window by its reconstruction error.
\end{itemize}

\emph{Additionally evaluated (not in the headline nine).} We also evaluate Deep SVDD \cite{ruff2018deep}---a one-class deep model whose MLP encoder maps normal data toward a fixed center, scoring a window by its squared distance to that center---and one-class SVM \cite{scholkopf2001estimating}. Deep SVDD produced no valid alarm on any FEMTO bearing at any training fraction (Section~\ref{sec:mintrain}) and is omitted from the headline set; its full numbers are retained, unchanged, in the result tables throughout (Tables~\ref{tab:imslead}, \ref{tab:xjtu}, \ref{tab:femto}, \ref{tab:ferrara}, \ref{tab:holm}, \ref{tab:farbudget}, \ref{tab:tradeoff}, \ref{tab:phrank}, and~\ref{tab:mintrain}). We keep it visible because it is the low-false-alarm operating point under the gated metric (Section~\ref{sec:tradeoff}, Table~\ref{tab:farbudget}). The deep sequence models carry an honest caveat: with on the order of 100 normal training windows they are data-starved. We include them both to provide a modern-baseline comparison and to determine empirically whether a finite commissioning-data requirement exists that makes deep reconstruction models competitive with SPC charts---and if so, to quantify it---a practically important question we address in Section~\ref{sec:mintrain}. We do not tune them per-dataset, the same fixed-protocol fairness applied to the SPC charts. On runs with fewer feature windows than the sequence length (short XJTU bearings, some coarse IMS factors) a sequence cannot be formed; we record that (run, detector) cell as an explicit N/A rather than dropping it silently. All non-conformal detectors alarm via a percentile threshold on training scores followed by a persistence filter, so the alarm semantics are identical across detectors and only the score differs.

\subsection{Conformal Calibration}
A conformal detector \cite{vovk2005algorithmic} wraps Isolation Forest. Nonconformity scores are computed on a calibration split of normal data; a test point alarms when its conformal $p$-value falls below the target $\alpha$. Under exchangeability of the calibration and test scores this guarantees $\text{FAR} \le \alpha$ in finite samples. We do not assume the guarantee: we validate it empirically by measuring $\text{FAR}_{\text{pre}}$ over the onset-defined normal region across a grid of $\alpha$, and we report the runs where it fails together with the reason.

\subsection{SCADA Sampling Sweep}
We coarsen the time grid by a factor $f \in \{1, 2, 5, 10, 20\}$ before feature extraction under two mechanisms:
\begin{itemize}
\item \emph{aggregate:} take the mean within each bin of width $f$, exactly as a SCADA historian stores a bin-averaged value;
\item \emph{decimate:} keep every $f$-th sample and discard the rest.
\end{itemize}
A naive sweep changes the window's wall-clock span as $f$ grows, which by itself alters lead time through a window-counting artifact rather than information loss. Our controlled sweep resamples at the feature level so that window content and alarm persistence are held constant and only the effective logging interval varies. This isolates the quantity of interest---does averaging lose fault-relevant information relative to decimation---from the confound.

\subsection{Statistical Protocol}
\label{sec:statproto}
Results are emitted in long format, one row per run $\times$ mode $\times$ factor $\times$ seed $\times$ method. Confidence intervals are obtained by bootstrapping runs, which are the unit of generalization (resampling seeds would understate uncertainty).

\emph{The run is the unit of inference.} A single physical bearing run produces one degradation trajectory; the five coarsening factors we apply to that run ($f \in \{1, 2, 5, 10, 20\}$) are not independent observations but repeated transformations of the same trajectory---statistical pseudoreplicates. Pooling them as $3 \times 5 = 15$ ``pairs'' per detector (as an earlier version of this work did) inflates the effective sample size fivefold and yields overconfident $p$-values. We therefore collapse each run to a single difference, the mean of (aggregate $-$ decimate) over factors, giving one observation per run: $n=3$ on IMS, $n=10$ on XJTU-SY, and $n=1$ on ONGC (a single asset, for which we report descriptive differences only and draw no inference).

On the resulting run-level differences we apply an exact two-sided sign test (binomial, $H_0 : P(\text{aggregate} > \text{decimate}) = \tfrac{1}{2}$). The sign test is assumption-light and the conservative choice at small $n$: it uses only the direction of each run's difference, making no normality or symmetry assumption that $n=3$ could not support. A parametric mixed-effects model with three random-effect levels would be over-fit and is deliberately not used. Note that the test cannot reach significance at $n=3$: the most extreme outcome, all three runs concordant, gives $p = 2(\tfrac{1}{2})^3 = 0.25$. We report this floor explicitly so that the IMS direction is read as a trend, not a significant effect. Runs whose mean (aggregate $-$ decimate) difference is exactly zero after factor-collapsing are excluded from the sign-test count; the $n_+$ and $n_-$ tallies in Tables~\ref{tab:xjtu} and \ref{tab:femto} count only non-zero-difference runs. Across the detector $\times$ dataset family of aggregate-vs-decimate tests (ONGC excluded, $n=1$) we control the family-wise error rate with the Holm--Bonferroni step-down procedure at $\alpha=0.05$, stating the family size $N$ and the uncorrected $\text{FWER} = 1 - (1-\alpha)^N$ alongside. We use a single fixed hyperparameter protocol throughout to avoid optimistic, per-condition tuning bias.


\section{Experimental Setup}
All experiments use a fixed random-seed protocol; stochastic detectors (Isolation Forest, the conformal wrapper) inject the seed explicitly so that runs are bit-reproducible. The feature pipeline caches extracted features to Parquet so that the sampling sweep, ablations, calibration, and trade-off analyses all consume identical features. The released environment is pinned to the exact package versions that produced the numbers reported here, and a test suite (32 tests) asserts the metric's ungameability properties, the leakage-free onset, and the conformal factory. The figures and tables in this paper are generated directly from the released result tables; no number in this manuscript is hand-edited.


\section{Results}
\label{sec:results}
We summarize the outcomes first, then support each in turn. We present IMS first for its raw-waveform mechanism access; XJTU-SY ($n=10$) is the primary inferential test.

\begin{itemize}
\item[(i)] \emph{The metric behaves.} Under the gated, onset-relative metric the detectors rank by genuine early warning, with the naive RMS-trend baseline correctly weakest (Table~\ref{tab:imslead}), and the metric's anchor---the degradation onset---is stable to its threshold (Table~\ref{tab:onsetsens}).
\item[(ii)] \emph{Aggregation does not destroy lead time---the central result.} On IMS it shows a consistent positive trend---the same sign in all three runs for every magnitude detector ($3\sigma$ $+18.4$~h median, IF $+12.2$~h, EWMA $+5.8$~h)---but at the honest unit of inference ($n=3$ runs) this is not statistically significant (Table~\ref{tab:imssweep}). On the real ONGC turbine the difference is about a minute ($n=1$ case study, Appendix~D.2, Table~\ref{tab:ongc}); on XJTU-SY ($n=10$, Table~\ref{tab:xjtu}), FEMTO/PRONOSTIA ($n=6$, Table~\ref{tab:femto}), and University of Ferrara ($n=6$, Table~\ref{tab:ferrara}) the effect is null, with Isolation Forest even reversing sign on FEMTO.
\item[(iii)] \emph{Nothing survives correction.} After Holm--Bonferroni across the $N=40$ detector $\times$ dataset family, no hypothesis is rejected (smallest adjusted $p = 1.00$, Table~\ref{tab:holm}): the defensible cross-dataset headline is non-destruction, not ``aggregation helps.''
\item[(iv)] \emph{The mechanism is generic smoothing.} The IMS advantage grows as noise rises (from $-0.35$~h to $+6.6$~h at 10~dB SNR) and is reproduced by a plain moving average or Kalman filter on the decimated stream (Tables~\ref{tab:noise}--\ref{tab:denoise})---so the cause is low-pass smoothing, not historian averaging specifically.
\item[(v)] \emph{Calibration and feature analysis.} The conformal wrapper attains its $\text{FAR} \le \alpha$ guarantee on the two runs that satisfy exchangeability and fails on the one with a late onset (Table~\ref{tab:conformal}); the trade-off and ablation analyses (Tables~\ref{tab:tradeoff}, \ref{tab:ablation}) confirm that simple time-domain amplitude/impulse features give the most valid early warning.
\end{itemize}
The remainder of this section develops each result.

\subsection{Lead Time and Confidence Intervals on IMS}
Table~\ref{tab:imslead} reports the mean raw (ungated) detection lead time across the three IMS runs at full resolution, with 95\% bootstrap confidence intervals obtained by resampling the $n=3$ runs. On raw lead alone the three deep reconstruction models top the table and Hotelling $T^2$ leads the classical detectors---but this ordering is exactly what the false-alarm budget overturns: the daggered detectors post their long leads only by flooding the pre-onset region, and have no valid operating point under $\tau=10\%$ (Table~\ref{tab:tradeoff}). Once the budget is enforced, $3\sigma$ and Deep SVDD are the deployable choices, a reordering we make precise as a formal comparison against the Saxena prognostic horizon in Section~\ref{sec:saxena}. The three deep reconstruction models post raw leads of 183--198~h but score $L=0$ on all three IMS runs---their valid-alarm fraction at that threshold is 0/3---because at the 97.5th-percentile operating threshold their pre-onset FAR ranges from 61\% on test~3 to 100\% on tests~1--2, far above the $\tau=10\%$ budget. Sweeping the alarm threshold changes this only once: across the percentiles of Table~\ref{tab:tradeoff}, TCN-AE and Transformer-AD attain no valid operating point at any level tested, while LSTM-AE attains one at the 99.5th percentile on test~3 alone---59.7~h of lead at 4.19\% pre-onset FAR, unchanged across ten random seeds---leaving its valid-alarm fraction at 1/3. The intervals are wide; this is the honest consequence of $n=3$ and is precisely why we add ONGC and the ten XJTU bearings next. Figure~\ref{fig:health} shows the leakage-free onset and failure markers on a representative run, with the shaded training region from which all baseline statistics are computed.

\begin{table}[t] \small
	\begin{center}
	\caption{\textbf{$\dagger$~Detectors with no valid operating point under the $\tau=10\%$ pre-onset-FAR budget (Table~\ref{tab:tradeoff}): their long raw leads are achieved only by flooding the pre-onset region with false alarms and are not usable warning time.} IMS mean detection lead time (hours) at full resolution, all ten evaluated detectors (the nine headline detectors plus Deep SVDD), mean with 95\% bootstrap CI (resampling $n=3$ runs). This is the raw, ungated lead.}
	\label{tab:imslead}
	\begin{tabular}{l r r}
		\hline \hline
		\textbf{Method} & \textbf{Lead time (h)} & \textbf{95\% CI} \\
		\hline
		\multicolumn{3}{l}{\emph{No valid operating point at $\tau=10\%$ (raw lead only)}} \\
		LSTM-AE$^\dagger$ & 214.9 & [65.5, 383.5] \\
		TCN-AE$^\dagger$ & 201.5 & [65.5, 343.5] \\
		Transformer-AD$^\dagger$ & 199.6 & [65.5, 337.7] \\
		EWMA$^\dagger$ & 78.3 & [27.2, 152.2] \\
		CUSUM$^\dagger$ & 77.2 & [26.3, 151.4] \\
		Isolation Forest$^\dagger$ & 69.8 & [28.0, 123.3] \\
		Hotelling $T^2{}^\dagger$ & 67.5 & [27.2, 113.3] \\
		\hline
		\multicolumn{3}{l}{\emph{Valid operating point exists---deployable under $\tau=10\%$}} \\
		$3\sigma$ & 58.0 & [22.2, 93.7] \\
		RMS-trend & 34.5 & [0.0, 93.7] \\
		Deep SVDD & 18.4 & [0.0, 45.5] \\
		\hline \hline
	\end{tabular}
	\end{center}
\end{table}

\begin{figure}[t]
\centering
\includegraphics[width=\columnwidth]{figure1}
\caption{Health-indicator trajectory for IMS Bearing test 2 with the leakage-free degradation onset and failure markers. The baseline mean and standard deviation use only the shaded training region, so the onset cannot see the test data.}
\label{fig:health}
\end{figure}

\subsection{Sensitivity of the Onset Definition}
Because the entire metric is anchored to the degradation onset $t_o$, its robustness matters. Table~\ref{tab:onsetsens} sweeps the threshold multiplier $k$ of Eq.~(\ref{eq:onset}) for the default health indicator (the RMS/kurtosis maximum, terminal-excursion rule). The onset position is stable: across $k \in \{3,4,5\}$ it moves by less than two percent of run span on every run, and the implied maximum lead time (failure minus onset) changes only modestly. The table also exposes an important structural fact about the first IMS test: its onset sits at ${\sim}98\%$ of the run span, i.e.\ only ${\sim}12.9$~h before failure. This run degrades abruptly, so its ``pre-onset'' region is almost the entire record and, critically, may already contain early degradation---a fact that returns to explain the conformal result below.

\begin{table}[t] \small
	\begin{center}
	\caption{Onset sensitivity on IMS (default RMS/kurtosis indicator, terminal-excursion rule). Onset position is reported as a percentage of run span; max lead is failure minus onset.}
	\label{tab:onsetsens}
	\begin{tabular}{l c r r}
		\hline \hline
		\textbf{Run} & \textbf{$k$} & \textbf{Onset (\% span)} & \textbf{Max lead (h)} \\
		\hline
		\multirow{3}{*}{1st\_test} & 3 & 98.4 & 12.9 \\
		 & 4 & 98.4 & 12.9 \\
		 & 5 & 98.5 & 12.7 \\
		\hline
		\multirow{3}{*}{2nd\_test} & 3 & 60.4 & 64.8 \\
		 & 4 & 61.9 & 62.5 \\
		 & 5 & 65.8 & 56.0 \\
		\hline
		\multirow{3}{*}{3rd\_test} & 3 & 94.5 & 59.3 \\
		 & 4 & 95.9 & 44.5 \\
		 & 5 & 96.1 & 42.2 \\
		\hline \hline
	\end{tabular}
	\end{center}
\end{table}

\emph{Short-bearing onset stability (XJTU and FEMTO).} A natural concern is that the fixed-sample onset parameters (persistence 5, gap tolerance 10) are a larger fraction of a sub-2-hour bearing than of a multi-week IMS run, so the onset could behave differently---or fail to fire---on the short campaigns. Repeating the $k \in \{3,4,5\}$ sweep on all sixteen short bearings shows it does not. Every bearing yields an onset at the default $k$, including the shortest runs in either campaign (XJTU Bearing2\_4 at 42~min / 0.68~h, FEMTO Bearing3\_1 at 1.43~h), so the persistence/gap filters do not preclude detection at short run length. The onset position is also stable to $k$: the median across-$k$ spread is 1.0 percentage points of run span on XJTU and 0.15 on FEMTO (versus 1.6 on IMS). We flag the two honest exceptions rather than hide them: XJTU Bearing1\_2 yields an onset at only one of the three $k$ values (its excursion is marginal), and FEMTO Bearing1\_1 has the widest spread (8.9 points)---both short-lived bearings whose marginal degradation the metric correctly treats as fragile. The stability claim therefore holds across all three run-to-failure datasets, with the fragility confined to exactly the bearings the per-bearing analysis already flags as hard.

\emph{Onset under a decoupled health indicator.} A sharper concern is circularity: the default onset indicator shares its RMS/kurtosis basis with the detectors' own features, so onset and detection are not fully independent (Section~\ref{sec:limitations}). We probe this directly by recomputing the onset under two alternatives that do not share the amplitude basis of the magnitude detectors---a kurtosis-only indicator (impulsiveness alone) and the first principal component (pca1) of the health channels (Table~\ref{tab:decoupled}). The result is honest and two-sided. The onset position is not invariant to the indicator: on the slow-degrading runs (tests 2 and 3) a kurtosis-only onset fires much later than the RMS-inclusive default (98.8\% vs 61.9\% of span on test 2), because on these bearings amplitude rises before impulsiveness, so an indicator that ignores RMS only reacts near failure. This means the absolute lead magnitudes and the per-run pre-onset region do depend on the indicator, and we do not claim otherwise. What is not affected is the paper's central conclusion: the aggregate-vs-decimate contrast is computed on raw lead time (Section~\ref{sec:statproto}), which contains no onset term, so non-destruction is invariant to the onset definition by construction. The circularity therefore bears on the absolute warning-time numbers and on which alarms count as valid---quantities we already report as run-specific---but not on the sign-of-difference headline. Test 1, the abrupt run, agrees across all three onset definitions (${\sim}98$--99\% of span).

\begin{table}[t] \small
	\begin{center}
	\caption{Onset position (\% of run span) and implied max lead (h) on IMS under three health-indicator definitions. The default rms\_kurt shares its amplitude basis with the magnitude detectors; kurt\_only and pca1 do not. The onset shifts materially on the slow-degrading tests 2 and 3, which we report honestly---but the non-destruction headline is computed on raw lead and is unaffected by this choice.}
	\label{tab:decoupled}
	\begin{tabular}{l c c c}
		\hline \hline
		\textbf{Run} & \textbf{rms\_kurt} & \textbf{kurt\_only} & \textbf{pca1} \\
		\hline
		1st\_test & 98.4\% [12.9] & 98.7\% [10.4] & 98.7\% [10.7] \\
		2nd\_test & 61.9\% [62.5] & 98.8\% [2.0] & 83.6\% [26.8] \\
		3rd\_test & 95.9\% [44.5] & 99.7\% [3.5] & 97.2\% [30.0] \\
		\hline \hline
	\end{tabular}
	\end{center}
\end{table}

\emph{Alarm-persistence sensitivity.} The alarm-persistence filter (default 3 windows) is a fixed protocol choice, so we sweep it over $\{1,3,5,10\}$ windows on IMS (Table~\ref{tab:persistence}). The two conclusions that matter are stable. First, the aggregate-vs-decimate direction stays non-negative at every persistence: the $3\sigma$ and EWMA charts show a sign-consistent positive run-level median at all four settings ($n_+/n_-=3/0$), and only the spike-robust Isolation Forest---the weakest carrier of the trend---drifts to a near-zero negative median at the longest persistence. Second, the $3\sigma$ valid-alarm fraction is 1.00 at every persistence, so the deployable-detector recommendation does not depend on this parameter.

\begin{table}[t] \small
	\begin{center}
	\caption{Alarm-persistence sensitivity on IMS: run-level median aggregate$-$decimate lead-time difference (h) and the $3\sigma$ valid-alarm fraction, at persistence $\in \{1,3,5,10\}$ windows. Non-destruction and the $3\sigma$ recommendation are stable throughout.}
	\label{tab:persistence}
	\begin{tabular}{l c c c c}
		\hline \hline
		\textbf{Persistence (windows)} & \textbf{1} & \textbf{3} & \textbf{5} & \textbf{10} \\
		\hline
		$3\sigma$ median $\Delta$ (h) & $+4.2$ & $+18.4$ & $+17.4$ & $+11.4$ \\
		EWMA median $\Delta$ (h) & $+5.2$ & $+5.8$ & $+4.8$ & $+3.0$ \\
		Iso.\ Forest median $\Delta$ (h) & $+10.4$ & $+12.2$ & $-0.8$ & $-0.2$ \\
		$3\sigma$ valid-alarm frac. & 1.00 & 1.00 & 1.00 & 1.00 \\
		\hline \hline
	\end{tabular}
	\end{center}
\end{table}

\emph{Missing-data (historian-gap) robustness.} Real historians drop records, so we inject random gaps---setting 5\% and 20\% of snapshot rows to missing (NaN, timestamp retained)---and re-measure lead time on IMS and XJTU (Table~\ref{tab:missing}). Warning time is robust to gaps of this magnitude: the SPC charts degrade negligibly ($3\sigma$ is unchanged on IMS through 20\% gaps; EWMA loses $<2$~h), and every detector's valid-alarm fraction is unchanged. The one visible sensitivity is Isolation Forest, whose IMS lead falls from 54.7 to 37.2~h at 20\% gaps---the same spike-sensitivity that makes it the weakest carrier of the aggregation trend. The finding is the missing-data analogue of non-destruction: dropping a fifth of the historian's records does not cost the control charts their warning time.

\begin{table}[t] \small
	\begin{center}
	\caption{Missing-data robustness: mean valid lead time (h) under random historian gaps of 0\%, 5\%, 20\% on IMS and XJTU-SY. Valid-alarm fractions (not shown) are unchanged across gap levels. The SPC charts are essentially gap-insensitive; only the spike-sensitive Isolation Forest degrades at 20\% on IMS.}
	\label{tab:missing}
	\begin{tabular}{l c c c | c c c}
		\hline \hline
		 & \multicolumn{3}{c|}{\textbf{IMS lead (h)}} & \multicolumn{3}{c}{\textbf{XJTU-SY lead (h)}} \\
		\textbf{Detector} & 0\% & 5\% & 20\% & 0\% & 5\% & 20\% \\
		\hline
		$3\sigma$ & 42.6 & 42.6 & 42.6 & 1.0 & 0.9 & 0.9 \\
		EWMA & 41.8 & 41.8 & 40.5 & 1.3 & 1.0 & 1.0 \\
		Iso.\ Forest & 54.7 & 53.4 & 37.2 & 1.2 & 1.2 & 1.1 \\
		\hline \hline
	\end{tabular}
	\end{center}
\end{table}

\subsection{Generalization to XJTU-SY ($n=10$ bearings)}
Ten run-to-failure bearings provide a far stronger test than IMS's three. Table~\ref{tab:xjtu} shows that the IMS ``aggregation helps'' trend does not replicate: at the run level (collapsing factors, $n=10$) the median difference is 0~h for six of the seven detectors ($-0.10$~h for Isolation Forest), the sign test is non-significant for all ($p \ge 0.18$), and no direction is consistent across bearings. Crucially, aggregation does not destroy lead time either---the medians are flat and the per-run differences are within a few minutes. Lead times are short in absolute terms (0--3.5~h) because XJTU bearings are short-lived (52--533~min); over the five window-magnitude detectors (the three SPC charts $3\sigma$/EWMA/Hotelling, Isolation Forest, and the RMS-trend floor---the same five tallied per bearing in Table~\ref{tab:perbearing}), 209/450 evaluations yield a valid alarm, with honest heterogeneity: the shortest-lived and the abrupt-failure bearings are hard or impossible to warn on under any mechanism. The deep sequence models are mostly N/A here---only the three longest-lived bearings furnish enough feature windows to form a length-30 sequence (76/90 cells N/A per deep model), an honest consequence of short runs rather than a tuning failure---and where they do run they too show a null aggregate-vs-decimate difference. Halving the sequence length to 15 windows does not change this picture: it makes only one additional short bearing usable and the deep reconstruction models still yield a valid alarm on at most one of the ten bearings, so trading sequence context for a marginally larger sample does not make them competitive with the SPC charts. Figure~\ref{fig:crossdataset} summarizes the cross-dataset picture in one view: a consistent but non-significant positive trend on IMS, and null effects on ONGC and XJTU.

\begin{table}[t] \small
	\begin{center}
	\caption{XJTU-SY ($n=10$): run-level sign test of aggregate vs decimate lead time, collapsing within-bearing sampling factors, for the seven non-sequence detectors. The three deep sequence models are $n=3$ on XJTU and are folded into Table~\ref{tab:holm}. No detector shows a significant or sign-consistent difference; Holm-adjusted $p$ in Table~\ref{tab:holm}.}
	\label{tab:xjtu}
	\begin{tabular}{l r c c r}
		\hline \hline
		\textbf{Method} & \textbf{Med.\ $\Delta$ (h)} & \textbf{$n_+/n_-$} & \textbf{Sign-con.} & \textbf{$p$} \\
		\hline
		$3\sigma$ & 0.00 & 3/1 & no & 0.625 \\
		EWMA & 0.00 & 2/2 & no & 1.000 \\
		CUSUM & 0.00 & 1/2 & no & 1.000 \\
		Hotelling $T^2$ & 0.00 & 2/0 & yes & 0.500 \\
		Iso.\ Forest & $-0.10$ & 2/7 & no & 0.180 \\
		Deep SVDD & 0.00 & 3/2 & no & 1.000 \\
		RMS-trend & 0.00 & 2/2 & no & 1.000 \\
		\hline \hline
	\end{tabular}
	\end{center}
\end{table}

\begin{figure*}[t]
\centering
\includegraphics[width=\textwidth]{fig_crossdataset}
\caption{Median run-level lead-time difference (aggregate $-$ decimate) per dataset and detector, annotated with the unit of inference ($n=3$ IMS, $n=10$ XJTU, $n=6$ FEMTO, $n=6$ Ferrara), plus the $n=1$ ONGC case study (Appendix~D.2). IMS shows a consistent positive trend; the inferential XJTU, FEMTO, and Ferrara differences are ${\approx}0$ on this scale. The honest cross-dataset headline is non-destruction. ONGC bars are hatched to mark the $n=1$ descriptive case study as separate from the four inferential datasets.}
\label{fig:crossdataset}
\end{figure*}

\subsection{Generalization to FEMTO/PRONOSTIA ($n=6$ bearings)}
The six FEMTO Learning\_set bearings provide a second, fully independent multi-bearing test, and---unlike XJTU, where most bearings are too short to form a length-30 sequence---all six are long enough that every detector, including the four deep reconstruction models, runs on the full $n=6$. Table~\ref{tab:femto} shows the same verdict as XJTU: at the run level the median aggregate$-$decimate difference is within $\pm9$~min for every detector, no detector is sign-consistent except Isolation Forest, and the smallest sign-test $p$ is 0.031 (Isolation Forest, raw)---which does not survive the Holm correction across the family (Table~\ref{tab:holm}). Notably, the one detector that is sign-consistent on FEMTO, Isolation Forest, trends negative (6/6 runs, median $-0.14$~h), the opposite sign to its positive IMS trend. This sign reversal across datasets is itself the cleanest statement of the finding: the direction of any residual aggregate-vs-decimate effect is dataset-specific noise, while the robust, replicated conclusion is non-destruction. Across all ten evaluated detectors (every one runs on the full $n=6$), 191/600 evaluations yield a valid alarm; as on XJTU, the short-lived bearings (${\sim}1.4$--7.8~h) keep absolute lead times small.

\begin{table}[t] \small
	\begin{center}
	\caption{FEMTO/PRONOSTIA ($n=6$): run-level sign test of aggregate vs decimate lead time. All ten evaluated detectors (nine headline plus Deep SVDD) run on the full $n=6$. No detector survives Holm correction (Table~\ref{tab:holm}); the only sign-consistent detector, Isolation Forest, trends negative---opposite to its IMS sign.}
	\label{tab:femto}
	\begin{tabular}{l r c c r}
		\hline \hline
		\textbf{Method} & \textbf{Med.\ $\Delta$ (h)} & \textbf{$n_+/n_-$} & \textbf{Sign-con.} & \textbf{$p$} \\
		\hline
		$3\sigma$ & $-0.10$ & 1/5 & no & 0.219 \\
		EWMA & 0.01 & 3/3 & no & 1.000 \\
		CUSUM & 0.00 & 3/3 & no & 1.000 \\
		Hotelling $T^2$ & 0.07 & 5/1 & no & 0.219 \\
		Iso.\ Forest & $-0.14$ & 0/6 & yes & 0.031 \\
		Deep SVDD & $-0.00$ & 3/3 & no & 1.000 \\
		LSTM-AE & $-0.05$ & 2/4 & no & 0.688 \\
		TCN-AE & 0.00 & 3/3 & no & 1.000 \\
		Transformer-AD & 0.04 & 5/1 & no & 0.219 \\
		RMS-trend & $-0.03$ & 2/4 & no & 0.688 \\
		\hline \hline
	\end{tabular}
	\end{center}
\end{table}

\subsection{Generalization to University of Ferrara ($n=6$ bearings)}
The University of Ferrara campaign \cite{arpa2024ferrara} adds a fifth, fully independent run-to-failure dataset and a second constant-condition multi-bearing rig beyond FEMTO. Table~\ref{tab:ferrara} gives the same verdict as XJTU-SY and FEMTO: at the run level the median aggregate$-$decimate difference is within $\pm5$~min for every detector, no detector is sign-consistent across the six bearings, and the smallest sign-test $p$ is 0.219 (Hotelling $T^2$)---not significant even before correction. Like the other short-lived campaigns, Ferrara's 0.7--6.8~h lifetimes mean it cannot confirm or deny the IMS positive trend; what it does add is a second independent replication of non-destruction on a modern (2024) rig. The lone few-minute negatives visible in the per-run lists (e.g.\ the $-1.0$ to $-1.5$~h first entries) are all on E1, whose degradation onset sits within ${\sim}0.1$~h of failure (an abrupt, near-coincident-onset bearing analogous to IMS test 1), not a systematic aggregation penalty.

\begin{table}[t] \small
	\begin{center}
	\caption{University of Ferrara 2024 ($n=6$): run-level sign test of aggregate vs decimate lead time. As on XJTU-SY and FEMTO the effect is null---no detector is sign-consistent and none survives Holm correction (Table~\ref{tab:holm}). $n_+/n_-$ count only non-zero-difference runs.}
	\label{tab:ferrara}
	\begin{tabular}{l r c c r}
		\hline \hline
		\textbf{Method} & \textbf{Med.\ $\Delta$ (h)} & \textbf{$n_+/n_-$} & \textbf{Sign-con.} & \textbf{$p$} \\
		\hline
		$3\sigma$ & $-0.02$ & 3/3 & no & 1.000 \\
		EWMA & 0.01 & 3/3 & no & 1.000 \\
		CUSUM & 0.01 & 4/2 & no & 0.688 \\
		Hotelling $T^2$ & 0.05 & 5/1 & no & 0.219 \\
		Iso.\ Forest & $-0.08$ & 2/4 & no & 0.688 \\
		Deep SVDD & $-0.01$ & 2/4 & no & 0.688 \\
		LSTM-AE & 0.00 & 4/2 & no & 0.688 \\
		TCN-AE & 0.00 & 4/2 & no & 0.688 \\
		Transformer-AD & 0.00 & 4/2 & no & 0.688 \\
		RMS-trend & $-0.00$ & 3/3 & no & 1.000 \\
		\hline \hline
	\end{tabular}
	\end{center}
\end{table}

\subsection{Aggregate vs Decimate on IMS: a Consistent Positive Trend}
Figure~\ref{fig:sweep} plots lead time against the effective logging interval for both mechanisms. Under aggregation, lead time remains high and stable as the interval coarsens; under decimation it degrades faster, with $3\sigma$ and Isolation Forest collapsing at coarse intervals (for example, $3\sigma$ falls from 58.0~h at full rate to 7.1~h at a $10\times$ decimation, while aggregation holds it near 59~h). The upturn at the very largest factor reflects the few-windows regime, where only one or two runs survive and the mean is no longer comparable; we annotate rather than hide it.

Table~\ref{tab:imssweep} gives the run-level test over all ten evaluated detectors. Honoring the run as the unit of inference (Section~\ref{sec:statproto}), the six magnitude-monitoring detectors all show a sign-consistent positive difference where the effect is non-trivial: $3\sigma$ (per-run $+27.1, +18.4, +15.6$~h; median $+18.4$), Isolation Forest ($+19.9, +12.2, +1.8$; median $+12.2$), EWMA ($+22.6, +5.0, +5.8$; median $+5.8$), CUSUM ($+22.6, +5.0, +5.4$; median $+5.4$), Hotelling $T^2$ ($+3.3, +12.3, +4.6$; median $+4.6$), and Deep SVDD ($+15.2, 0.0, +1.8$; two positive, one tie). The three deep reconstruction models behave oppositely but at negligible magnitude: LSTM-AE, TCN, and Transformer-AD each show small negative per-run differences (medians $-1.4, -0.4, -1.4$~h), consistent with the smoothing mechanism---a reconstruction model reconstructs the smoothed signal about as well either way, so it does not inherit the threshold-crossing advantage that the variance-sensitive charts get. The RMS-trend baseline is mixed. The direction is therefore a consistent positive trend for the chart/distance detectors---in no run does aggregation shorten their lead time---but with only $n=3$ independent runs the exact sign test attains its floor of $p=0.25$ (Section~\ref{sec:statproto}), so it is not statistically significant and does not survive Holm correction. We make no significance claim; the defensible reading is non-destruction with a consistently positive direction for the magnitude-monitoring detectors. The candidate mechanism is that bin-averaging acts as a low-pass filter suppressing single-window noise in the health statistic, so the alarm threshold is crossed earlier and held more persistently than under decimation. We test this smoothing hypothesis directly in Appendix~D.1.

\begin{figure*}[t]
\centering
\includegraphics[width=\textwidth]{fig_sweep}
\caption{Detection lead time vs effective SCADA logging interval on IMS (controlled sweep, mean across runs). Aggregation (left) preserves lead time as the interval coarsens far better than decimation (right), where the variance-sensitive detectors collapse at coarse rates.}
\label{fig:sweep}
\end{figure*}

\begin{table}[t] \small
	\begin{center}
	\caption{IMS controlled sweep: run-level difference of aggregate vs decimate lead time, collapsing the five within-run sampling factors to one mean difference per run (the unit of inference; $n=3$). The exact two-sided sign test floors at $p=0.25$ at $n=3$. Holm-adjusted $p$ in Table~\ref{tab:holm}.}
	\label{tab:imssweep}
	\begin{tabular}{l l r c r}
		\hline \hline
		\textbf{Method} & \textbf{Run diffs (h)} & \textbf{Med.} & \textbf{S-c.} & \textbf{$p$} \\
		\hline
		$3\sigma$ & $+27.1,+18.4,+15.6$ & $+18.4$ & yes & 0.25 \\
		Iso.\ Forest & $+19.9,+12.2,+1.8$ & $+12.2$ & yes & 0.25 \\
		EWMA & $+22.6,+5.0,+5.8$ & $+5.8$ & yes & 0.25 \\
		CUSUM & $+22.6,+5.0,+5.4$ & $+5.4$ & yes & 0.25 \\
		Hotelling $T^2$ & $+3.3,+12.3,+4.6$ & $+4.6$ & yes & 0.25 \\
		Deep SVDD & $+15.2,0.0,+1.8$ & $+1.8$ & 2+,tie & 0.50 \\
		RMS-trend & $-19.4,0.0,+2.8$ & 0.0 & no & 1.00 \\
		LSTM-AE & $-1.4,-0.4,-1.7$ & $-1.4$ & yes ($-$) & 0.25 \\
		TCN-AE & $-1.4,-0.4,-0.2$ & $-0.4$ & yes ($-$) & 0.25 \\
		Transformer-AD & $-1.4,-0.4,-2.5$ & $-1.4$ & yes ($-$) & 0.25 \\
		\hline \hline
	\end{tabular}
	\end{center}
\end{table}

\subsection{Multiple-Comparison Correction}
The aggregate-vs-decimate question is asked once per detector on each inferential dataset, so the honest family of hypotheses is the $10 \times 4 = 40$ detector $\times$ dataset tests over IMS, XJTU-SY, FEMTO, and Ferrara (ONGC, $n=1$, is excluded). The uncorrected family-wise error rate at $\alpha=0.05$ is $1 - (1-0.05)^{40} \approx 0.87$, which is exactly why a correction is required. We apply Holm--Bonferroni to the forty run-level sign-test $p$-values (Table~\ref{tab:holm}). No hypothesis survives (0 of 40): the smallest raw $p$-value is 0.031 (Isolation Forest on FEMTO), and after Holm correction across the $N=40$ family its adjusted $p$ is 1.00, still far above $\alpha$. The conclusion is unambiguous and conservative: after honoring the run structure and correcting for multiple comparisons, there is no statistically significant aggregate-vs-decimate effect on any dataset. The IMS sign-consistency remains a real, reportable trend; it is simply not, and cannot be at $n=3$, a significant one. We note that the family is heterogeneous in power: the three deep sequence models contribute $n=3$ XJTU tests rather than $n=10$ (marked $\ddagger$ in Table~\ref{tab:holm}), and those three cells rest on the longest-lived---hence most warnable---bearings. Because every uncorrected $p$ in the family already exceeds 0.03, this heterogeneity does not change any rejection decision; we flag it so the $N=40$ count is not read as forty equally-powered tests.

\begin{table}[t] \small
	\begin{center}
	\caption{Holm--Bonferroni correction across the $N=40$ detector $\times$ dataset family (IMS $n=3$, XJTU-SY $n=10$, FEMTO $n=6$, Ferrara $n=6$; ONGC excluded, $n=1$). No hypothesis is rejected at $\alpha=0.05$ (0 of 40 survive). $\ddagger$: the three deep sequence models need a length-30 window sequence, so those XJTU tests are $n=3$. The $N=40$ family is the full evaluated set of ten detectors (nine headline plus the additionally-evaluated Deep SVDD), retained exactly as originally computed.}
	\label{tab:holm}
	\begin{tabular}{l l r r c}
		\hline \hline
		\textbf{Dataset} & \textbf{Method} & \textbf{Raw $p$} & \textbf{Holm $p$} & \textbf{Rej.?} \\
		\hline
		IMS & $3\sigma$ & 0.250 & 1.000 & no \\
		IMS & EWMA & 0.250 & 1.000 & no \\
		IMS & CUSUM & 0.250 & 1.000 & no \\
		IMS & Hotelling $T^2$ & 0.250 & 1.000 & no \\
		IMS & Iso.\ Forest & 0.250 & 1.000 & no \\
		IMS & Deep SVDD & 0.500 & 1.000 & no \\
		IMS & LSTM-AE & 0.250 & 1.000 & no \\
		IMS & TCN-AE & 0.250 & 1.000 & no \\
		IMS & Transformer-AD & 0.250 & 1.000 & no \\
		IMS & RMS-trend & 1.000 & 1.000 & no \\
		XJTU & $3\sigma$ & 0.625 & 1.000 & no \\
		XJTU & Iso.\ Forest & 0.180 & 1.000 & no \\
		XJTU & Hotelling $T^2$ & 0.500 & 1.000 & no \\
		XJTU & CUSUM & 1.000 & 1.000 & no \\
		XJTU & EWMA & 1.000 & 1.000 & no \\
		XJTU & Deep SVDD & 1.000 & 1.000 & no \\
		XJTU & LSTM-AE$^\ddagger$ & 1.000 & 1.000 & no \\
		XJTU & TCN-AE$^\ddagger$ & 1.000 & 1.000 & no \\
		XJTU & Transf.-AD$^\ddagger$ & 1.000 & 1.000 & no \\
		XJTU & RMS-trend & 1.000 & 1.000 & no \\
		FEMTO & $3\sigma$ & 0.219 & 1.000 & no \\
		FEMTO & EWMA & 1.000 & 1.000 & no \\
		FEMTO & CUSUM & 1.000 & 1.000 & no \\
		FEMTO & Hotelling $T^2$ & 0.219 & 1.000 & no \\
		FEMTO & Iso.\ Forest & 0.031 & 1.000 & no \\
		FEMTO & Deep SVDD & 1.000 & 1.000 & no \\
		FEMTO & LSTM-AE & 0.688 & 1.000 & no \\
		FEMTO & TCN-AE & 1.000 & 1.000 & no \\
		FEMTO & Transf.-AD & 0.219 & 1.000 & no \\
		FEMTO & RMS-trend & 0.688 & 1.000 & no \\
		Ferrara & $3\sigma$ & 1.000 & 1.000 & no \\
		Ferrara & EWMA & 1.000 & 1.000 & no \\
		Ferrara & CUSUM & 0.688 & 1.000 & no \\
		Ferrara & Hotelling $T^2$ & 0.219 & 1.000 & no \\
		Ferrara & Iso.\ Forest & 0.688 & 1.000 & no \\
		Ferrara & Deep SVDD & 0.688 & 1.000 & no \\
		Ferrara & LSTM-AE & 0.688 & 1.000 & no \\
		Ferrara & TCN-AE & 0.688 & 1.000 & no \\
		Ferrara & Transf.-AD & 0.688 & 1.000 & no \\
		Ferrara & RMS-trend & 1.000 & 1.000 & no \\
		\hline \hline
	\end{tabular}
	\end{center}
\end{table}

\emph{Per-bearing heterogeneity.} Table~\ref{tab:perbearing} breaks the XJTU result down by bearing at full resolution (aggregate mechanism, valid alarms counted over the five detectors). We report the mean lead over the valid detectors as the realistic single-deployed-detector expectation, alongside the oracle ``best lead'' (the maximum over five, which presumes foreknowledge of which detector fires earliest); the two differ only where detectors disagree (e.g.\ Bearing2\_2: 1.00~h mean vs.\ 1.08~h oracle). The aggregate 209/450 valid-alarm figure quoted above is not uniform: it is the sum of a strongly heterogeneous per-bearing picture. Long-degrading bearings such as Bearing2\_3 (533~min life) and Bearing1\_3 (158~min) are warnable with multi-hour or near-hour lead, whereas the two shortest-lived, abrupt-failure bearings---Bearing2\_4 (42~min) and Bearing1\_1 (123~min with a sudden end)---yield zero valid alarms for any detector (2/10 bearings are thus unwarnable by every method), because their degradation onset and failure are nearly coincident and no method can manufacture warning time that the physics does not provide. This is the honest face of the metric: it refuses to credit warning time on bearings that genuinely give none, which is exactly the behavior an operator should want from an evaluation.

\begin{table}[t] \small
	\begin{center}
	\caption{XJTU-SY per-bearing summary at full resolution (aggregate). ``Valid/5'' counts the detectors (of five) that produce a valid alarm; ``Mean valid'' is the mean lead over those detectors; ``Best'' is the oracle maximum over the five. Two bearings are warned by no detector (0/5).}
	\label{tab:perbearing}
	\begin{tabular}{l c r c r r}
		\hline \hline
		\textbf{Bearing} & \textbf{Cd.} & \textbf{Life} & \textbf{V/5} & \textbf{Mean} & \textbf{Best} \\
		\hline
		Bearing1\_1 & 1 & 123 & 0 & -- & 0.00 \\
		Bearing1\_2 & 1 & 161 & 1 & 0.67 & 0.67 \\
		Bearing1\_3 & 1 & 158 & 4 & 1.03 & 1.03 \\
		Bearing1\_4 & 1 & 122 & 2 & 0.68 & 0.68 \\
		Bearing1\_5 & 1 & 52 & 3 & 0.35 & 0.35 \\
		Bearing2\_1 & 2 & 491 & 3 & 0.67 & 0.67 \\
		Bearing2\_2 & 2 & 161 & 2 & 1.00 & 1.08 \\
		Bearing2\_3 & 2 & 533 & 1 & 3.45 & 3.45 \\
		Bearing2\_4 & 2 & 42 & 0 & -- & 0.00 \\
		Bearing2\_5 & 2 & 339 & 4 & 2.30 & 2.30 \\
		\hline
		\multicolumn{4}{l}{Total (full-res, aggregate)} & \multicolumn{2}{r}{20/50} \\
		\multicolumn{4}{l}{Bearings with no valid detector} & \multicolumn{2}{r}{2/10} \\
		\hline \hline
	\end{tabular}
	\end{center}
\end{table}

\subsection{Conformal Calibration}
Table~\ref{tab:conformal} reports the empirical pre-onset false-alarm rate of the conformal Isolation Forest against the target $\alpha$ on each IMS run, and Figure~\ref{fig:conf_ims} plots it. The distribution-free guarantee $\text{FAR} \le \alpha$ holds on test 2 (0\% through $\alpha=0.10$) and is essentially satisfied on test 3 (0.9\% at $\alpha=0.01$, 4.4\% at $\alpha=0.05$---so it holds through $\alpha=0.05$---while it slightly exceeds the bound at $\alpha=0.10$, where the empirical FAR is 13\%). It is violated on test 1 for a principled reason that Table~\ref{tab:onsetsens} already foretold: that run's onset fires only ${\sim}12$~h before failure, so its ``pre-onset'' region spans almost the whole record and already contains degradation; exchangeability between the calibration and ``normal'' scores therefore fails, and the empirical FAR sits at 0.74--0.85 regardless of $\alpha$. We report this run rather than discarding it, because it delimits exactly when the conformal guarantee is and is not usable in prognostics: it requires a genuinely stationary normal region, which an abrupt late-onset failure does not provide. The XJTU calibration (Figure~\ref{fig:conf_xjtu}) tells the same story at scale: bearings with a substantial stationary pre-onset region (e.g.\ Bearing2\_3, calibrated through $\alpha=0.05$) honor the bound, while those with only one or two pre-onset windows---an artifact of the short XJTU lifetimes---cannot satisfy a finite-sample guarantee and are honestly flagged uncalibrated rather than silently passed. The conformal machinery thus doubles as a per-asset usability test: it reports whether enough stationary normal data exists to control the false-alarm rate at all.

\emph{FEMTO and ONGC calibration.} The same per-asset usability test on FEMTO (Figure~\ref{fig:conf_femto}) and the ONGC turbine (Figure~\ref{fig:conf_ongc}) confirms the pattern and its honest limits. On FEMTO, no bearing honors the bound at $\alpha=0.05$ (empirical pre-onset FAR 0.08--0.56; pooled 0.27), because the short 1.4--7.8~h lifetimes leave a pre-onset region that is neither long nor stationary---the same reason the short XJTU bearings and IMS test 1 fail. The two bearings with the most stationary normal region (Bearing2\_2 at 0.08, Bearing1\_2 at 0.10) come closest, while the abrupt Bearing1\_1 (0.56) violates it badly---exactly the diagnostic the conformal machinery is meant to provide. On the real ONGC historian stream the single-asset calibration tracks the diagonal much more closely (empirical FAR 0.016, 0.030, 0.096 at $\alpha=0.01, 0.02, 0.05$), sitting a modest 1.5--2$\times$ above the target---consistent with a genuine, mostly-stationary operating record whose normal region only mildly violates exchangeability. The conformal guarantee is therefore usable where a stationary normal region exists (ONGC, the longer IMS/XJTU bearings) and correctly self-reports as invalid where it does not (short FEMTO/XJTU bearings, abrupt onsets).

\begin{table}[t] \small
	\begin{center}
	\caption{Conformal Isolation Forest on IMS: empirical pre-onset FAR vs target $\alpha$. \checkmark\ marks $\text{FAR} \le \alpha$ (guarantee satisfied).}
	\label{tab:conformal}
	\begin{tabular}{l c c c c}
		\hline \hline
		\textbf{Run} & \textbf{$\alpha{=}0.01$} & \textbf{$\alpha{=}0.05$} & \textbf{$\alpha{=}0.10$} & \textbf{$\alpha{=}0.20$} \\
		\hline
		1st\_test & 0.74 & 0.82 & 0.85 & 0.86 \\
		2nd\_test & 0.00\,\checkmark & 0.00\,\checkmark & 0.00\,\checkmark & 0.25 \\
		3rd\_test & 0.01\,\checkmark & 0.04\,\checkmark & 0.13 & 0.30 \\
		\hline \hline
	\end{tabular}
	\end{center}
\end{table}

\begin{figure}[t]
\centering
\includegraphics[width=\columnwidth]{fig_conf_ims}
\caption{Conformal calibration on IMS: empirical pre-onset FAR vs target $\alpha$. Points on or below the diagonal satisfy the distribution-free guarantee. Test 1 violates it because its near-failure onset leaves no genuinely stationary normal region.}
\label{fig:conf_ims}
\end{figure}

\begin{figure}[t]
\centering
\includegraphics[width=\columnwidth]{fig_conf_xjtu}
\caption{Conformal calibration on the XJTU-SY bearings. Bearings with a substantial stationary pre-onset region honor the bound; those with only one or two pre-onset windows cannot support a finite-sample guarantee and are flagged uncalibrated.}
\label{fig:conf_xjtu}
\end{figure}

\begin{figure}[t]
\centering
\includegraphics[width=\columnwidth]{fig_conf_femto}
\caption{Conformal calibration on the six FEMTO/PRONOSTIA bearings. No bearing honors the bound at $\alpha=0.05$; the short, non-stationary pre-onset regions of these 1.4--7.8~h bearings mean the finite-sample guarantee cannot hold, and the machinery flags them uncalibrated.}
\label{fig:conf_femto}
\end{figure}

\subsection{Lead-Time--vs--False-Alarm Trade-off}
\label{sec:tradeoff}
Sweeping the alarm-threshold percentile traces a prognostics analog of the ROC curve (Figure~\ref{fig:tradeoff_ims}, Table~\ref{tab:tradeoff}): tightening the threshold lowers both pre-onset FAR and lead time together, letting an operator select an operating point for a stated false-alarm budget. EWMA and Isolation Forest dominate the upper-left of the curve---long lead at moderate FAR---while the naive RMS-trend baseline collapses to near-zero lead as soon as the threshold is tightened. At a strict 99.5th-percentile threshold, for instance, $3\sigma$ retains 56.3~h of lead at 8.9\% pre-onset FAR, whereas RMS-trend has collapsed to 1.3~h. This converts the single gated number into a usable operating-point selector. Note that many entries in Table~\ref{tab:tradeoff} sit above the $\tau=10\%$ false-alarm budget; under the gated metric those operating points are invalid and we mark them with a dagger. This has a concrete consequence for detector choice: EWMA attains the highest raw lead in Table~\ref{tab:tradeoff} but its pre-onset FAR exceeds the $\tau=10\%$ budget at every threshold tested (all three EWMA entries are daggered), so under the gated metric it has no valid operating point on IMS. CUSUM behaves the same way: its pre-onset FAR sits at 19.0--19.4\% across the swept thresholds, never dipping below the budget, so it too has no valid operating point---the two memory-based control charts share this liability, and neither should be recommended under a 10\% budget. Deep SVDD occupies the opposite corner: it stays comfortably within budget ($\le 0.6\%$ FAR) but only by firing late, buying that low false-alarm rate with a short $\le 14.8$~h lead. The $3\sigma$ chart, by contrast, reaches a valid point (56.3~h of lead at 8.9\% FAR at the 99.5th percentile) and is therefore the recommended chart under a 10\% budget---a ranking that the raw-lead column of Table~\ref{tab:imslead}, where EWMA leads, would by itself obscure. The same trade-off computed on the ten XJTU bearings (Figure~\ref{fig:tradeoff_xjtu}) has the same qualitative shape at the shorter absolute lead times those bearings allow, confirming that the operating-point selection generalizes beyond the laboratory rig.

\emph{Sensitivity to the false-alarm budget.} The budget $\tau$ is a policy input, so we report how the detector recommendation moves with it rather than fixing $\tau=10\%$ silently (Table~\ref{tab:farbudget}). The recommendation is not $\tau$-invariant, and we state this plainly. At $\tau=0.20$ the $3\sigma$ chart gives the longest valid lead (77.0~h), now joined by CUSUM, Hotelling $T^2$, and Isolation Forest as their pre-onset FAR falls within the looser budget. At $\tau=0.10$ the $3\sigma$ chart remains the best valid chart (56.3~h). But at a strict $\tau=0.05$ no control chart qualifies at all---only Deep SVDD (14.8~h) and the RMS-trend floor (34.7~h) stay within a 5\% budget---so an operator with a very low nuisance-alarm tolerance should prefer the low-FAR Deep SVDD operating point over any chart. The $3\sigma$ recommendation therefore holds for $\tau \ge 0.10$ on IMS and inverts toward the low-FAR detectors below it. Crucially, the cross-dataset non-destruction headline is untouched by this choice: it is computed on raw (ungated) lead-time differences (Section~\ref{sec:statproto}) and does not depend on $\tau$ at all.

\emph{Where does $\tau=10\%$ come from? A cost-ratio argument.} We do not treat $\tau$ as arbitrary. Consider a simple expected-cost model with two site parameters: $c_I$, the cost of responding to one nuisance alarm (a wasted inspection), and $c_F$, the cost of an unplanned bearing failure (downtime plus secondary shaft/housing damage). Over a bearing's pre-onset life of $k$ windows a detector at pre-onset FAR $f$ issues ${\approx} kf$ nuisance inspections, and a valid alarm with adequate lead converts the failure into a planned repair, saving $c_F$. An operator should tolerate false alarms up to the budget at which their cumulative cost equals the warning's value, $k\tau c_I \approx c_F$, i.e.
\begin{equation}
\tau^* \approx \frac{c_F}{k\,c_I} = \frac{\rho}{k}, \qquad \rho \equiv c_F/c_I.
\label{eq:costratio}
\end{equation}
For rolling-element bearings an unplanned failure typically costs $\rho = 10^2$--$10^4$ times a routine inspection, and the pre-onset window count is of order $k \sim 10^3$ on the runs studied here, so Eq.~(\ref{eq:costratio}) places the break-even budget at $\tau^* \sim 0.1$--1 (Table~\ref{tab:taugrid}): our $\tau=0.10$ is the conservative end of that range and is justified whenever failures are at least ${\sim}100\times$ costlier than inspections. The strict $\tau=0.05$ column of Table~\ref{tab:farbudget} corresponds to the opposite regime---inspections that are themselves disruptive (e.g.\ requiring a shutdown), so $\rho$ is small---exactly where the low-FAR Deep SVDD point becomes the economical choice. We stress this is an order-of-magnitude anchor, not a site-specific optimisation: the true optimum requires an operator's actual $c_I$ and $c_F$, which we do not have, and Eq.~(\ref{eq:costratio}) is offered so a practitioner can substitute their own.

\begin{figure}[t]
\centering
\includegraphics[width=\columnwidth]{fig_tradeoff_ims}
\caption{Lead-time vs pre-onset-FAR trade-off on IMS (threshold percentile swept, mean across runs). Up-and-to-the-left is better. EWMA and Isolation Forest dominate; the RMS-trend baseline collapses at low FAR.}
\label{fig:tradeoff_ims}
\end{figure}

\begin{figure}[t]
\centering
\includegraphics[width=\columnwidth]{fig_tradeoff_xjtu}
\caption{Lead-time vs pre-onset-FAR trade-off on the ten XJTU-SY bearings. The curve has the same qualitative shape as IMS at the shorter absolute lead times XJTU's short-lived bearings permit, showing the operating-point selection generalizes.}
\label{fig:tradeoff_xjtu}
\end{figure}

\begin{figure}[t]
\centering
\includegraphics[width=\columnwidth]{fig_farbudget}
\caption{Best valid lead time per detector as a function of the false-alarm budget $\tau$ on IMS, plotted directly from the operating points of Table~\ref{tab:farbudget}. Each line is one detector; a value of 0 means no operating point within budget. The recommended detector changes with $\tau$: $3\sigma$ wins at $\tau \ge 0.10$, while at a strict $\tau=0.05$ only the low-FAR Deep SVDD and the RMS-trend floor remain valid.}
\label{fig:farbudget_fig}
\end{figure}

\begin{table}[t] \small
	\begin{center}
	\caption{FAR-budget sensitivity on IMS: best valid lead time (h) per detector at each false-alarm budget $\tau$ (0.0 = no operating point within budget). The $3\sigma$ recommendation holds at $\tau=0.10$ and 0.20 but not at a strict $\tau=0.05$. The non-destruction conclusion is computed on raw lead and is unaffected by $\tau$.}
	\label{tab:farbudget}
	\begin{tabular}{l c c c}
		\hline \hline
		\textbf{Detector} & \textbf{$\tau{=}0.05$} & \textbf{$\tau{=}0.10$} & \textbf{$\tau{=}0.20$} \\
		\hline
		$3\sigma$ & 0.0 & 56.3 & 77.0 \\
		EWMA & 0.0 & 0.0 & 0.0 \\
		CUSUM & 0.0 & 0.0 & 65.0 \\
		Hotelling $T^2$ & 0.0 & 0.0 & 73.9 \\
		Iso.\ Forest & 0.0 & 0.0 & 53.4 \\
		Deep SVDD & 14.8 & 14.8 & 14.8 \\
		RMS-trend & 34.7 & 34.7 & 34.7 \\
		\hline \hline
	\end{tabular}
	\end{center}
\end{table}

\begin{table}[t] \small
	\begin{center}
	\caption{The break-even budget $\tau^* = \rho/k$ from Eq.~(\ref{eq:costratio}) (clipped to $[0,1]$), over failure-to-inspection cost ratios $\rho \in \{10^2, 10^3, 10^4\}$ and pre-onset window counts $k \in \{10^2, 10^3, 10^4\}$. The chosen $\tau=0.10$ sits inside this range across the mid-to-high cost ratios typical of rolling-element bearings. Every entry is computed directly from Eq.~(\ref{eq:costratio}); no experimental quantity enters.}
	\label{tab:taugrid}
	\begin{tabular}{l c c c}
		\hline \hline
		 & \textbf{$\rho{=}10^2$} & \textbf{$\rho{=}10^3$} & \textbf{$\rho{=}10^4$} \\
		\hline
		$k{=}10^2$ & 1.00 & 1.00 & 1.00 \\
		$k{=}10^3$ & 0.10 & 1.00 & 1.00 \\
		$k{=}10^4$ & 0.01 & 0.10 & 1.00 \\
		\hline \hline
	\end{tabular}
	\end{center}
\end{table}

\begin{table}[t] \small
	\begin{center}
	\caption{IMS lead-time--vs--false-alarm trade-off (mean across runs) at three threshold percentiles. Up-and-to-the-left (high lead, low FAR) is better. $\dagger$ marks operating points whose pre-onset FAR exceeds the $\tau=10\%$ budget and are therefore invalid under the gated metric.}
	\label{tab:tradeoff}
	\begin{tabular}{l r r r r r r}
		\hline \hline
		 & \multicolumn{2}{c}{\textbf{95th}} & \multicolumn{2}{c}{\textbf{99th}} & \multicolumn{2}{c}{\textbf{99.5th}} \\
		\textbf{Method} & Ld & FAR & Ld & FAR & Ld & FAR \\
		\hline
		$3\sigma$ & 77.0 & 18.1$\dagger$ & 60.7 & 11.7$\dagger$ & 56.3 & 8.9 \\
		EWMA & 89.1 & 22.1$\dagger$ & 78.3 & 21.0$\dagger$ & 78.3 & 21.0$\dagger$ \\
		CUSUM & 65.0 & 19.4$\dagger$ & 64.2 & 19.0$\dagger$ & 64.2 & 19.0$\dagger$ \\
		Hot.\ $T^2$ & 176.9 & 28.1$\dagger$ & 82.7 & 20.8$\dagger$ & 73.9 & 19.0$\dagger$ \\
		Iso.\ For. & 174.7 & 27.1$\dagger$ & 75.8 & 20.6$\dagger$ & 53.4 & 19.8$\dagger$ \\
		Deep SVDD & 14.8 & 0.6 & 3.2 & 0.0 & 2.7 & 0.0 \\
		RMS-trend & 34.7 & 4.6 & 1.3 & 0.0 & 1.3 & 0.0 \\
		\hline \hline
	\end{tabular}
	\end{center}
\end{table}

\subsection{Feature-Group Ablation}
Table~\ref{tab:ablation} ablates feature families on IMS for Isolation Forest, reporting mean lead time, mean pre-onset FAR, and the valid-alarm fraction across the three runs. The pattern is instructive and somewhat against expectation. Spectral and envelope groups yield the longest raw lead times (99.4~h for envelope) but the lowest validity: they fire early in part because they fire indiscriminately, accruing 23--29\% pre-onset FAR. The rms\_only group, by contrast, yields the highest valid-alarm fraction (1.00) at a near-zero pre-onset FAR (0.2\%), at the cost of a shorter raw lead. In other words, once the false-alarm budget is enforced, the simple time-domain amplitude statistics, not the elaborate frequency-domain features, carry the usable prognostic signal on IMS.

To check that this recommendation is not specific to Isolation Forest, we repeat the ablation for the four SPC charts and Deep SVDD (Table~\ref{tab:ablation_general}, valid-alarm fraction by group). The conclusion holds across every non-sequence detector: time-domain amplitude features (rms\_only or time\_only) give each detector its highest or tied-highest validity, and admitting spectral, envelope, or the full feature set never improves validity for any of them. For the memory-based control charts it sharpens: rms\_only is the unique group at which $3\sigma$ and EWMA reach full validity (1.00 at 0\% pre-onset FAR), and spectral/envelope/all collapse all four SPC charts to 0.33 by driving their pre-onset FAR to 42--63\%---the variance-sensitive charts react violently to the noisy extra dimensions. Deep SVDD is the instructive exception, and it fails for the opposite reason: its pre-onset FAR stays near 0\% under every feature set (it rarely false-alarms), yet spectral features still cut its validity to 0.33 (and to 0.00 for spectral\_only) by delaying or suppressing the alarm rather than inflating the FAR---the frequency-domain dimensions hurt the one-class deep model through missed detections, not false alarms. Either way the practical recommendation is detector-general on IMS: under a false-alarm budget, restrict every detector to time-domain amplitude features; frequency-domain features must be admitted with care, most of all for the SPC charts.

\emph{Robustness of the rms\_only recommendation to the logging rate.} Because the recommendation is meant to guide SCADA deployments, we repeat the feature-group ablation at every coarsening factor (Table~\ref{tab:coarsen_ablation}). At native and $2\times$ resolution rms\_only is the uniquely best group (valid fraction 1.00), matching the full-resolution ablation and the informational argument that bin-averaging preserves the amplitude statistic---the mean of a bin-averaged signal's RMS tracks the RMS of the original, so rms\_only carries essentially the same signal after aggregation. This validation holds cleanly at native and $2\times$ resolution ($f=1,2$), where rms\_only is uniquely best and each IMS run contributes enough windows to distinguish feature groups. At $f \ge 5$ only a handful of windows survive per run, so valid-alarm fractions move in steps of 1/3 and the ranking reflects few-window noise rather than stable feature-group differences---at $f=20$, for instance, time\_only recovers to 0.89 while rms\_only sits at 0.67, but this reversal rests on single-window differences per run and cannot be read as a stable recommendation. The defensible reading is therefore: rms\_only is the safest feature set at logging rates that retain enough windows to decide ($f \le 2$), and no feature group displaces it at native resolution. Whether the rms\_only advantage holds at the coarsest SCADA rates is not resolvable at $n=3$ and is left as an open empirical question.

\begin{table}[t] \small
	\begin{center}
	\caption{IMS feature-group ablation (Isolation Forest): mean lead time, pre-onset FAR, and valid-alarm fraction across the three runs. Validity, not raw lead, is the operationally meaningful column.}
	\label{tab:ablation}
	\begin{tabular}{l c r r r}
		\hline \hline
		\textbf{Feature group} & \textbf{$d$} & \textbf{Lead (h)} & \textbf{FAR (\%)} & \textbf{Valid} \\
		\hline
		rms\_only & 8 & 32.0 & 0.2 & 1.00 \\
		kurt\_only & 5 & 60.2 & 16.2 & 0.67 \\
		time\_only & 20 & 63.4 & 18.8 & 0.67 \\
		spectral\_only & 30 & 70.6 & 23.4 & 0.67 \\
		envelope\_only & 16 & 99.4 & 29.3 & 0.67 \\
		time+spectral & 50 & 70.9 & 20.6 & 0.67 \\
		all & 50 & 70.9 & 20.6 & 0.67 \\
		\hline \hline
	\end{tabular}
	\end{center}
\end{table}

\begin{table}[t] \small
	\begin{center}
	\caption{Feature-group valid-alarm fraction by coarsening factor on IMS ($3\sigma$, EWMA, Isolation Forest; mean over the three runs and the three detectors). rms\_only is uniquely best at native and near-native resolution ($f\le2$); at $f \ge 5$ the surviving-window count is small and the fractions are coarse (1/3 steps).}
	\label{tab:coarsen_ablation}
	\begin{tabular}{l c c c c c}
		\hline \hline
		\textbf{Feature group} & \textbf{$f{=}1$} & \textbf{$f{=}2$} & \textbf{$f{=}5$} & \textbf{$f{=}10$} & \textbf{$f{=}20$} \\
		\hline
		rms\_only & 1.00 & 1.00 & 0.56 & 0.44 & 0.67 \\
		time\_only & 0.67 & 0.67 & 0.78 & 0.56 & 0.89 \\
		spectral\_only & 0.44 & 0.56 & 0.67 & 0.56 & 0.67 \\
		all & 0.44 & 0.56 & 0.67 & 0.56 & 0.67 \\
		\hline \hline
	\end{tabular}
	\end{center}
\end{table}

\begin{table}[t] \small
	\begin{center}
	\caption{Feature-group ablation generalized beyond Isolation Forest: valid-alarm fraction (mean over the three IMS runs) by group. Time-domain amplitude features are best for every detector; for the SPC charts rms\_only is uniquely ($3\sigma$, EWMA) or jointly best, and spectral/envelope/all collapse them to 0.33.}
	\label{tab:ablation_general}
	\begin{tabular}{l c c c c c c}
		\hline \hline
		\textbf{Group} & \textbf{$3\sigma$} & \textbf{EW} & \textbf{CU} & \textbf{$T^2$} & \textbf{DS} & \textbf{IF} \\
		\hline
		rms\_only & 1.00 & 1.00 & 0.67 & 0.67 & 0.67 & 1.00 \\
		kurt\_only & 0.67 & 0.67 & 0.33 & 0.67 & 0.33 & 0.67 \\
		time\_only & 0.67 & 0.67 & 0.67 & 0.67 & 0.33 & 0.67 \\
		spectral\_only & 0.33 & 0.33 & 0.33 & 0.33 & 0.00 & 0.67 \\
		envelope\_only & 0.33 & 0.33 & 0.33 & 0.33 & 0.33 & 0.67 \\
		time+spectral & 0.33 & 0.33 & 0.33 & 0.33 & 0.33 & 0.67 \\
		all & 0.33 & 0.33 & 0.33 & 0.33 & 0.33 & 0.67 \\
		\hline \hline
	\end{tabular}
	\end{center}
\end{table}

\subsection{Spectral vs Time-Domain Features, per Run}
A dedicated time-vs-spectral comparison (Table~\ref{tab:spectral}) confirms the ablation at the level of individual runs and alarms. Adding spectral features to the time-domain set does not improve valid detection on IMS, and on test 2 it actively harms the control charts: the $3\sigma$, EWMA, and Hotelling detectors jump to 100\% pre-onset FAR (invalid) when spectral features are included, while the same detectors are valid with time-only features. Isolation Forest, which is less sensitive to a few noisy dimensions, remains valid in both settings. The first test is unwarnable for every configuration because of its near-failure onset, consistent with Table~\ref{tab:onsetsens} and the conformal result. The practical lesson is that frequency-domain features must be admitted with care under a false-alarm budget: their extra sensitivity cuts both ways.

\begin{table}[t] \small
	\begin{center}
	\caption{IMS spectral ablation, per run and detector (all six non-sequence detectors): valid alarm (\checkmark) or invalid ($\times$) under time-only vs time+spectral features. Spectral features inflate pre-onset FAR for the control charts on test 2, flipping each from valid to invalid; Isolation Forest is robust and stays valid. Test 1 is unwarnable for every configuration.}
	\label{tab:spectral}
	\begin{tabular}{l l c c}
		\hline \hline
		\textbf{Run} & \textbf{Detector} & \textbf{time-only} & \textbf{time+spec.} \\
		\hline
		\multirow{6}{*}{1st\_test} & $3\sigma$ & $\times$ & $\times$ \\
		 & EWMA & $\times$ & $\times$ \\
		 & CUSUM & $\times$ & $\times$ \\
		 & Hotelling $T^2$ & $\times$ & $\times$ \\
		 & Deep SVDD & $\times$ & $\times$ \\
		 & Iso.\ Forest & $\times$ & $\times$ \\
		\hline
		\multirow{6}{*}{2nd\_test} & $3\sigma$ & \checkmark & $\times$ \\
		 & EWMA & \checkmark & $\times$ \\
		 & CUSUM & \checkmark & $\times$ \\
		 & Hotelling $T^2$ & \checkmark & $\times$ \\
		 & Deep SVDD & $\times$ & $\times$ \\
		 & Iso.\ Forest & \checkmark & \checkmark \\
		\hline
		\multirow{6}{*}{3rd\_test} & $3\sigma$ & \checkmark & \checkmark \\
		 & EWMA & \checkmark & \checkmark \\
		 & CUSUM & \checkmark & \checkmark \\
		 & Hotelling $T^2$ & \checkmark & \checkmark \\
		 & Deep SVDD & \checkmark & \checkmark \\
		 & Iso.\ Forest & \checkmark & \checkmark \\
		\hline \hline
	\end{tabular}
	\end{center}
\end{table}

\subsection{Comparison with the Saxena Prognostic Horizon}
\label{sec:saxena}
If the false-alarm gate is what separates a usable warning from a manufactured one, then a metric without that gate should reward detectors our metric rejects. We test this against the most widely used ungated warning-time metric---the Saxena prognostic horizon (PH) \cite{saxena2008metrics,saxena2010metrics}---and show that on IMS it makes a different, and operationally wrong, detector recommendation.

The PH is the interval between a detector's first alarm and failure, $\text{PH} = \max(0, (t_f - \text{FAT})/3600)$~hours, with no false-alarm constraint. It is exactly our lead time $L$ before the validity gate (Eq.~\ref{eq:valid}) is applied; that one missing conjunct is decisive.

\emph{PH is maximized by the latch-on detector.} A detector that alarms from the first sample earns the entire record as prognostic horizon. Table~\ref{tab:latchon} evaluates this constant-on detector on the three IMS runs: its PH is the full run length (164--1073~h), yet it floods the whole pre-onset region ($\text{FAR}_{\text{pre}} \approx 1$), so the gated metric scores it invalid on every run ($L=0$). The oracle that fires just after the true onset $t_o^+$ is the mirror image---zero pre-onset false alarms and the maximum achievable lead. Under PH the latch-on detector is the clear winner; under $L$ it is correctly scored worthless.

\emph{PH and $L$ disagree on the real ranking.} The latch-on case is a worst-case existence proof; the more damaging point is that PH and $L$ disagree on deployable detectors evaluated at their own best operating points. Table~\ref{tab:phrank} ranks the seven non-sequence detectors on IMS by their best raw lead (PH) and by their best lead at an operating point within the $\tau=10\%$ budget ($L$). The two rankings are nearly inverted at the top. The four detectors PH ranks highest---Hotelling $T^2$ (176.9~h), Isolation Forest (174.7~h), EWMA (89.1~h), and CUSUM (65.0~h)---each have no valid operating point under a 10\% false-alarm budget, so their gated lead is $L=0$. The detector PH ranks only fourth, $3\sigma$ (77.0~h raw), is the one that actually clears the budget (56.3~h of lead at 8.9\% pre-onset FAR) and is the correct recommendation; Deep SVDD, ranked last of the seven by PH, rises to a valid 14.8~h. The three deep reconstruction models (omitted from the sweep) would sit above all of these by PH---raw leads of 200--215~h (Table~\ref{tab:imslead})---while alarming through 48--100\% of every pre-onset region, i.e.\ $L=0$: the most extreme illustration that PH rewards precisely the detectors a false-alarm-budgeted operator must avoid.

\emph{Implication.} Any warning-time metric that does not gate on $\text{FAR}_{\text{pre}}$ can be driven to its maximum by loosening the alarm threshold---in the limit, the constant-on detector. Our metric closes this by construction; the Saxena PH does not. The distinction is not theoretical: on IMS it changes the recommended detector from the ones PH ranks first by raw lead to $3\sigma$, the only classical chart with a valid operating point at a 10\% budget.

\begin{table}[t] \small
	\begin{center}
	\caption{Latch-on detector under the Saxena prognostic horizon vs.\ the gated metric $L$, on the three IMS runs. A constant-on detector earns the full run length as PH but floods the pre-onset region, so it is scored invalid ($L=0$). The oracle fires just after the true onset $t_o^+$: zero pre-onset false alarms and the maximum achievable lead.}
	\label{tab:latchon}
	\begin{tabular}{l l r c r c}
		\hline \hline
		\textbf{Detector} & \textbf{Run} & \textbf{PH} & \textbf{FAR} & \textbf{$L$} & \textbf{Val?} \\
		\hline
		\multirow{3}{*}{Latch-on} & 1st & 827.6 & ${\approx}1$ & 0.0 & no \\
		 & 2nd & 163.8 & ${\approx}1$ & 0.0 & no \\
		 & 3rd & 1073.3 & ${\approx}1$ & 0.0 & no \\
		\hline
		\multirow{3}{*}{Oracle} & 1st & 12.9 & 0.00 & 12.9 & yes \\
		 & 2nd & 62.5 & 0.00 & 62.5 & yes \\
		 & 3rd & 44.5 & 0.00 & 44.5 & yes \\
		\hline \hline
	\end{tabular}
	\end{center}
\end{table}

\begin{table}[t] \small
	\begin{center}
	\caption{Detector rankings on IMS: Saxena prognostic horizon (PH, ungated best raw lead) vs.\ the gated metric $L$ (best lead at an operating point within the $\tau=10\%$ budget). The rankings invert at the top: the detectors PH ranks best have no valid operating point, while $3\sigma$---PH's fourth---is the deployable winner. $^*$RMS-trend is valid only at the loosest threshold and on one of three runs.}
	\label{tab:phrank}
	\begin{tabular}{l r c r c c}
		\hline \hline
		\textbf{Method} & \textbf{PH} & \textbf{Rk} & \textbf{$L$} & \textbf{$L$ Rk} & \textbf{Valid?} \\
		\hline
		Hotelling $T^2$ & 176.9 & 1 & 0.0 & 4 & no \\
		Iso.\ Forest & 174.7 & 2 & 0.0 & 4 & no \\
		EWMA & 89.1 & 3 & 0.0 & 4 & no \\
		$3\sigma$ & 77.0 & 4 & 56.3 & 1 & yes \\
		CUSUM & 65.0 & 5 & 0.0 & 4 & no \\
		RMS-trend & 34.7 & 6 & 34.7 & 2 & yes$^*$ \\
		Deep SVDD & 14.8 & 7 & 14.8 & 3 & yes \\
		\hline \hline
	\end{tabular}
	\end{center}
\end{table}

\subsection{Minimum Training Data for Deep Anomaly Detectors}
\label{sec:mintrain}
A practitioner commissioning condition monitoring at a new asset must decide how much normal run-in history to collect before a deep model can be trusted over a simple control chart. We quantify this on FEMTO---the one dataset where all six bearings are long enough for every detector, including the seq-len-30 deep models, to run on the full $n=6$---by sweeping the training fraction $T \in \{0.20, \dots, 0.60\}$ (calibration fixed at 0.10, test $= 1-T-0.10$) with no per-$T$ retuning. At no training fraction did any bearing fail to form a length-30 sequence, so every cell is a genuine evaluation rather than an N/A.

The result (Table~\ref{tab:mintrain}, Figure~\ref{fig:mintrain}) is unambiguous and, we think, more useful than a crossover point would have been: there is no crossover. The four magnitude control charts ($3\sigma$, EWMA, CUSUM, Isolation Forest) reach their highest valid-alarm fraction (0.67--0.83 over the six bearings) with as little as $T=0.20$ of run length and are stable or only mildly lower thereafter. The three deep reconstruction models sit far below them at every training fraction (LSTM-AE and Transformer-AD 0.33--0.50, TCN-AE 0.17--0.33) and do not close the gap as $T$ grows---more commissioning data does not make them competitive in this regime. Deep SVDD produces no valid alarm on any FEMTO bearing at any $T$. The failure mode is the same as on IMS test 1: Deep SVDD fires very late in the run---its threshold requires substantial encoder-distance growth from the trained center before alarming---so on short FEMTO bearings (1.4--7.8~h) the first alarm arrives after or coincident with failure on every bearing regardless of training fraction, giving $L=0$ uniformly. This is a bearing-lifetime constraint, not a training-data constraint: the model is not data-starved in the usual sense but structurally unable to alarm early enough on runs this short. The valid fraction of the charts if anything edges down slightly from $T=0.20$ to $T=0.30$, an artifact of the shrinking test window on these short bearings rather than a data-hunger effect.

The practical recommendation is therefore sharper than ``collect more data'': for short-lived assets like FEMTO and XJTU bearings, SPC charts should be preferred regardless of how much run-in history is available, because the deep reconstruction models do not overtake them anywhere in the 0.20--0.60 training range. Whether a crossover appears on genuinely long-running assets---where a deep model could see thousands rather than dozens of normal windows---is the central open question. The closest publicly available single-bearing example is the KAIST bearing with a $\sim$43-hour lifecycle ($n=1$, insufficient for inference as noted in Section~\ref{sec:limitations}). Research groups with access to multi-run, long-runway industrial historian archives---gas turbines, compressors, paper machines---can apply the released pipeline directly; loader templates for all public dataset formats used in this study are included in the repository at \url{https://github.com/Ansh-Goyal01/Scada-Leadtime-Benchmark} to minimize integration effort.

\begin{figure}[t]
\centering
\includegraphics[width=\columnwidth]{fig_mintrain}
\caption{Valid-alarm fraction vs.\ training fraction on FEMTO ($n=6$). SPC charts (blue) sit at the top and are flat from $T=0.20$; the deep reconstruction models (red) and Deep SVDD lie below the SPC baseline (dashed) at every training fraction and never reach it---no crossover.}
\label{fig:mintrain}
\end{figure}

\begin{table}[t] \small
	\begin{center}
	\caption{Valid-alarm fraction vs.\ training fraction on FEMTO ($n=6$ bearings; no per-$T$ retuning). The SPC charts are stable and best from $T=0.20$ onward; the deep reconstruction models and Deep SVDD stay below them at every training fraction---no crossover in the 0.20--0.60 range. ``SPC mean'' is the mean over $3\sigma$/EWMA/CUSUM/IF.}
	\label{tab:mintrain}
	\begin{tabular}{l c c c c c}
		\hline \hline
		\textbf{$T$} & \textbf{$3\sigma$} & \textbf{EW} & \textbf{IF} & \textbf{DS} & \textbf{SPC} \\
		\hline
		0.20 & 0.83 & 0.83 & 0.83 & 0.00 & 0.83 \\
		0.30 & 0.83 & 0.67 & 0.67 & 0.00 & 0.71 \\
		0.40 & 0.83 & 0.67 & 0.67 & 0.00 & 0.71 \\
		0.50 & 0.67 & 0.67 & 0.67 & 0.00 & 0.67 \\
		0.60 & 0.67 & 0.67 & 0.67 & 0.00 & 0.67 \\
		\hline \hline
	\end{tabular}
	\end{center}
\end{table}


\FloatBarrier
\section{Discussion}

\subsection{What Changed the Story}
Four methodological corrections, applied together, change the conclusion of the sampling study. First, a correct failure label for the third IMS test (which had left a fifth of the data after the nominal failure) turns a guaranteed miss into a valid detection. Second, an onset-gated metric that defeats the latch-on exploit prevents a constant-on detector from manufacturing lead time. Third, a $p < n$, channel-invariant feature space removes the ill-posedness of a channel-count-dependent scheme. Fourth, a statistical test on the genuine unit of inference---the run, not the within-run sampling factor---replaces an eyeballed curve and a pseudoreplicated paired test. With all four in place, the original claim that SCADA averaging destroys lead time is refuted on every dataset tested, and the previously-reported significant positive effect is corrected to a consistent but non-significant trend.

\subsection{Why the IMS Trend Appears but Is Neutral Elsewhere}
The IMS trend has a clean candidate mechanism: bin-averaging is a low-pass filter that removes single-window noise spikes, so the variance-sensitive control charts ($3\sigma$, EWMA) cross and hold their thresholds earlier under aggregation than under decimation, which keeps the spikes. This is consistent with the trend concentrating in those detectors and being weakest for Isolation Forest, which is robust to isolated spikes by design. The IMS trend is reproduced by an ordinary moving-average or Kalman filter on the decimated stream (Appendix~D.1), indicating a generic low-pass-smoothing mechanism rather than historian averaging specifically. On ONGC the data is already smoothed at the 10~s historian rate, so further aggregation adds little; on XJTU the bearings are short-lived and many fail abruptly, so there is little quiet runway over which averaging could help. The honest synthesis is therefore not ``averaging helps'' but non-destruction: across very different regimes, coarse aggregation never costs warning time, and where the signal is noisy and the runway long it shows a consistent tendency to stabilize the health statistic.

\emph{The single-asset turbine case study.} The one real historian record in our study---an ONGC gas-turbine low-pressure-compressor bearing logged at a 10~s SCADA rate and ending in a genuine operator shutdown---has every detector alarm roughly 35~h ahead of the labeled failure, with an aggregate-vs-decimate difference of at most about a minute for every detector. Because it is a single asset ($n=1$) we draw no inference from it and exclude it from the inferential family; we report it in full as a case study in Appendix~D.2. It is nonetheless the cleanest illustration of the practical question---it uses no simulated coarsening---and is consistent with non-destruction.

\subsection{Practical Implications}
For practitioners the message is reassuring and actionable. Storing bin-averaged vibration at SCADA rates---the default behavior of most historians---does not inherently sacrifice bearing-fault warning time relative to keeping decimated raw samples, and may even improve it for control-chart detectors on noisy long-runway data. The conformal results add a usable, distribution-free false-alarm guarantee wherever the normal region is genuinely stationary, and the trade-off curve turns the gated metric into an operating-point selector for a stated false-alarm budget. The ablations caution that more features are not better under a false-alarm budget: time-domain amplitude and impulsiveness statistics carry the usable signal on IMS, and spectral features must be admitted with care. Operators whose cost structure penalizes false alarms heavily---for instance, process plants where each alarm triggers a mandatory manual inspection---have a complementary operating point available. Deep SVDD consistently stays within budget at $\le 0.6\%$ pre-onset FAR across the swept threshold levels on IMS (Table~\ref{tab:tradeoff}), achieving 14.8~h of advance warning with near-zero false alarms. Where $3\sigma$ maximizes warning time within the 10\% budget, Deep SVDD minimizes nuisance alarms for operators requiring near-zero tolerance. Neither should be used on IMS test 1, which the conformal machinery correctly flags as uncalibrated (Table~\ref{tab:conformal}).

\subsection{Summary of Key Outcomes}
In one place, the numbers that matter. (1) Non-destruction holds on all four inferential datasets (and the single-asset case study): SCADA bin-averaging never costs lead time relative to decimation. (2) On IMS, aggregation gives a positive but non-significant trend---same sign in all three runs for every magnitude detector ($3\sigma$ $+18.4$~h median, IF $+12.2$~h, EWMA $+5.8$~h, CUSUM $+5.4$~h, Hotelling $T^2$ $+4.6$~h), but the exact sign test floors at $p=0.25$ at $n=3$. (3) Elsewhere the effect is null: the single-asset ONGC case study differs by $\le 1$~min, and XJTU-SY ($n=10$), FEMTO ($n=6$), and Ferrara ($n=6$) by $\le 8$~min median, with Isolation Forest reversing sign on FEMTO. (4) After Holm correction across the $N=40$ detector $\times$ dataset family, no hypothesis survives (smallest adjusted $p=1.00$). (5) Mechanism: the IMS advantage grows from $-0.35$~h to $+6.6$~h as SNR falls to 10~dB and is matched by a plain moving-average or Kalman filter on the decimated stream---so the cause is generic low-pass smoothing, not historian averaging specifically. The honest, replicated headline is therefore non-destruction, not ``aggregation helps.''


\section{Limitations and Threats to Validity}
\label{sec:limitations}
We state the limits of the evidence plainly, grouped by the kind of validity each concern bears on.

\emph{External validity (generalization).} The IMS ``aggregation helps'' tendency is a sign-consistent but non-significant trend ($n=3$; Table~\ref{tab:imssweep}) and does not replicate on the larger XJTU (Table~\ref{tab:xjtu}) or FEMTO (Table~\ref{tab:femto}) sets---on FEMTO the only sign-consistent detector even trends in the opposite direction; we therefore claim only non-destruction as the cross-dataset headline, not a positive effect. ONGC is a single asset ($n=1$), reported as a descriptive case study with no inference (Table~\ref{tab:ongc}). The datasets, while deliberately diverse, are all rolling-element bearings; we do not claim the result transfers to gearboxes, blades, or other failure modes without further testing.

\emph{Residual gaming surfaces of the metric.} We claim the metric defeats the specific latch-on (constant-on) exploit, not that it is unconditionally ungameable, and we name two surfaces we do not close. (a) A detector could in principle align its alarm threshold to fire just after the onset time $t_o^+$; because our onset is computed only from training-region statistics, this cannot use test-region information, but an adversary with knowledge of the onset rule could still tune toward it. (b) The onset health indicator shares its RMS/kurtosis basis with several of the detectors' own features, so onset and detection are not fully independent---a circularity that could, in principle, favor detectors aligned to the onset statistic. We acknowledge both rather than claim to have eliminated them; closing them (e.g.\ an onset signal disjoint from the detector features) is future work.

\emph{Construct validity (does the metric measure warning time?).} The onset definition, although leakage-free and stable to its threshold multiplier across $k \in \{3,4,5\}$ (Table~\ref{tab:onsetsens}), is not the ground-truth defect-initiation time, which is unobservable in these datasets. A late onset---as on IMS test 1, where degradation and failure are nearly coincident---inflates the pre-onset region, can mark a genuine early detection invalid, and breaks the conformal exchangeability assumption (Table~\ref{tab:conformal}). The per-bearing XJTU breakdown (Table~\ref{tab:perbearing}) shows the same effect: bearings whose physics provides essentially no warning runway score zero valid alarms regardless of detector, which is correct behavior but means the valid-alarm fraction is partly a property of the bearings, not only of the methods.

\emph{Internal validity (is the comparison fair?).} We use a single fixed hyperparameter protocol throughout to avoid optimistic per-condition tuning, and we expose threshold sensitivity through the full trade-off sweep (Table~\ref{tab:tradeoff}) rather than reporting a single hand-picked operating point. The controlled sampling sweep holds window content and persistence constant so that the aggregate-vs-decimate contrast isolates information loss; the few-windows regime at the largest coarsening factor is annotated rather than hidden. Stochastic detectors inject their seed explicitly so results are bit-reproducible.

\emph{Statistical-conclusion validity.} Confidence intervals are wide on IMS ($n=3$); we resample runs, the genuine unit of generalization, rather than seeds, which would understate uncertainty. Crucially, the five sampling factors applied within a run are pseudoreplicates of one degradation trajectory, not independent observations: an earlier pooled 15-pair Wilcoxon test treated them as independent and thereby overstated significance. We collapse to one difference per run and apply an exact sign test, which at $n=3$ cannot fall below $p=0.25$; after Holm correction across the forty-test family nothing is significant (Table~\ref{tab:holm}). The IMS result is therefore a directional trend to be read together with the XJTU and FEMTO nulls, never as a significant effect. A parametric mixed-effects model is not used: with three random-effect levels it would be over-fit, and the assumption-light sign test is the defensible choice at this sample size. A formal power analysis makes the minimum-dataset requirement concrete: to detect the observed IMS effect (Cohen's $d \approx 0.7$--1.2 across the magnitude detectors) at $\alpha=0.05$ and power 0.80, a paired test needs $n \approx 6$ (at $d=1.2$) to $n \approx 17$ (at $d=0.7$) independent run-to-failure bearings; equivalently, the exact sign test we use cannot reach significance until at least six runs are sign-concordant ($2(\tfrac{1}{2})^6=0.031$). The three IMS runs are thus a factor of two to five short of the sample size a significance claim would require---the quantitative content of the ``long-runway, many-bearing campaign'' the future work calls for. Long-runway bearing datasets with multiple runs are rare in the public domain; the closest available single-bearing example is a KAIST bearing with a ${\sim}43$-hour lifecycle ($n=1$, insufficient for inference). Research groups with access to long-running industrial historian archives---gas turbines, compressors, paper machines---can directly apply our released pipeline to test this hypothesis; we provide loader templates for all public dataset formats used here to minimize integration effort.


\section{Conclusion}
We presented a lead-time-centric evaluation protocol for bearing anomaly detection---one that defeats the latch-on exploit through an onset-gated, false-alarm-budgeted metric---and used it to test whether SCADA-rate logging destroys detection lead time. Across four inferential run-to-failure datasets---XJTU-SY ($n=10$), FEMTO/PRONOSTIA ($n=6$), University of Ferrara ($n=6$), and NASA IMS ($n=3$)---together with a real single-asset turbine SCADA case study (Appendix~D.2), coarse aggregation does not destroy lead time---and on IMS shows a consistent, sign-concordant positive trend for the variance-sensitive control charts (not statistically significant at $n=3$, and not surviving multiple-comparison correction across the $N=40$ family)---while decimation is never better. The contribution is a reproducible methodology and an honest, counter-intuitive empirical result: non-destruction, not a guaranteed positive effect. Future work will learn the degradation onset jointly with the detector so that the metric's anchor is independent of a hand-set health indicator (closing the circularity noted in Section~\ref{sec:limitations}), and extend the conformal guarantee to the non-exchangeable early-degradation regime in which it currently fails. A third direction concerns the minimum training-data requirement for deep reconstruction models (Section~\ref{sec:mintrain}): verifying whether the crossover training fraction generalizes to longer-running assets and different machinery types is a natural extension. We release the full harness, detectors, tests, and the exact tables and figures of this paper to support replication and extension.


\FloatBarrier
\section*{Acknowledgments}
The author thanks the operators and data engineers at ONGC who collected and maintained the gas-turbine SCADA record used in the single-asset case study. The NASA/IMS, XJTU-SY, FEMTO/PRONOSTIA, and University of Ferrara run-to-failure datasets are publicly released by their respective providers and are gratefully acknowledged. All experiments were conducted on personal computing resources; no external funding was received for this work.


\section*{Data and Code Availability}
All code---detectors, the leakage-free onset, the gated metric, the sampling sweep, the conformal wrapper, the run-level statistical harness, and the noise/denoiser mechanism test---is released with a pinned environment (requirements.lock.txt, Python 3.13) and a test suite that asserts the metric's resistance to the latch-on exploit, the leakage-free onset, and the run-level sign-test/Holm machinery. Every table and figure in this paper is generated directly from the released result files; no number is hand-edited. Repository: \url{https://github.com/Ansh-Goyal01/Scada-Leadtime-Benchmark}. A versioned snapshot is archived on Zenodo at the concept DOI \url{https://doi.org/10.5281/zenodo.20719076}. The IMS (NASA), XJTU-SY, and FEMTO/PRONOSTIA run-to-failure datasets are publicly available from their original providers; the repository ships the loaders and the derived per-bearing feature caches needed to reproduce every result. The ONGC gas-turbine SCADA record is proprietary and covered by a non-disclosure agreement, so its raw data cannot be redistributed; the derived health-indicator series and the figures/tables it produces are included so the $n=1$ case study remains inspectable, and no headline statistical claim depends on it.


% Bibliography
% ---------------------------------------------------------------------------------
\bibliographystyle{apacite}
\PHMbibliography{ijphm}
% ---------------------------------------------------------------------------------


\FloatBarrier
\section*{Appendix A: Reproducibility Settings}
Table~\ref{tab:hyperparams} lists the fixed hyperparameters used throughout, taken verbatim from the released config.py. A single protocol is applied to every detector, dataset, run, and sampling factor---no per-condition tuning. Table~\ref{tab:compute} reports the wall-clock training and per-window inference cost of each detector on the largest IMS run ($X_{\text{train}}=631$, $X_{\text{test}}=506$ windows, single CPU core, median of five repeats). Training is a one-off offline step; inference is negligible for every detector---at most ${\sim}54$~$\mu$s per window (LSTM-AE), five orders of magnitude below the fastest historian polling interval in our study (10~s). Streaming condition monitoring at native SCADA rates is therefore not compute-bound for any of the ten evaluated detectors.

\begin{table}[t] \small
	\begin{center}
	\caption{Fixed hyperparameter protocol (from config.py). The onset persistence/gap tolerances are in samples of the health-indicator series. The pre-onset FAR is measured over the full pre-onset window $t < t_o$ (Eq.~\ref{eq:valid}).}
	\label{tab:hyperparams}
	\begin{tabular}{l l}
		\hline \hline
		\textbf{Setting} & \textbf{Value} \\
		\hline
		Feature window / overlap & 10 samples / 50\% \\
		Feature schema & channel-invariant, 49-dim \\
		Train-fit feature cap & top 50 (stratified) \\
		Train / calib.\ / test & 0.50 / 0.10 / 0.40 \\
		Pre-onset FAR region & full window $t < t_o$ \\
		Onset health indicator & rms\_kurt \\
		Onset method / persistence & terminal excursion / 5 \\
		Onset $k\sigma$ / gap tolerance & $4\sigma$ / 10 samples \\
		Alarm threshold / persistence & 97.5th pct / 3 windows \\
		FAR budget $\tau$ & 0.10 \\
		Isolation Forest & 200 trees, contam.\ auto \\
		Hotelling $T^2$ & PCA rank 10, $\alpha=0.01$ \\
		EWMA / CUSUM & $\lambda{=}0.2, k{=}3$ / $k{=}0.5, h{=}5$ \\
		Deep models & seq-len 30, 40--50 epochs \\
		Transformer & $d_{\text{model}}{=}32$, 2 heads, 2 layers \\
		FEMTO channels & horiz.\ + vert.\ accel. \\
		Conformal calib.\ fraction & 0.10 \\
		Bootstrap resamples & runs, $B{=}2000$ \\
		Random seed & 42 \\
		Environment & Py 3.13, requirements.lock.txt \\
		\hline \hline
	\end{tabular}
	\end{center}
\end{table}

\begin{table}[t] \small
	\begin{center}
	\caption{Per-detector compute cost on the largest IMS run (631 train / 506 test windows, single CPU core, median of five repeats). Even the slowest model's inference is $\ll$ the 10~s SCADA polling interval.}
	\label{tab:compute}
	\begin{tabular}{l l r}
		\hline \hline
		\textbf{Detector} & \textbf{Train} & \textbf{Inf.\ ($\mu$s/win)} \\
		\hline
		$3\sigma$ & $<1$ ms & $<0.1$ \\
		EWMA & $<1$ ms & 0.2 \\
		CUSUM & $<1$ ms & 0.5 \\
		Hotelling $T^2$ & 4 ms & 0.4 \\
		Isolation Forest & 0.20 s & 24 \\
		Deep SVDD & 0.41 s & 1.4 \\
		LSTM-AE & 7.8 s & 54 \\
		TCN-AE & 10.5 s & 47 \\
		Transformer-AD & 10.3 s & 38 \\
		\hline \hline
	\end{tabular}
	\end{center}
\end{table}

\FloatBarrier
\section*{Appendix B: The 49-Dimensional Invariant Feature Space}
For full reproducibility we give the exact channel-invariant feature construction. Each snapshot's raw waveform first yields six per-channel time-domain statistics (Table~\ref{tab:basestats}). Each family is then made channel-count-invariant by aggregating across the physical channels with two operators---the channel mean and the channel max (the worst-channel max is the classic fault indicator)---and each aggregated series is summarised over the rolling window (10 samples, 50\% overlap) by four descriptors: the window mean, standard deviation, maximum, and least-squares linear-trend slope. This yields $6 \times 2 \times 4 = 48$ features; one cross-channel RMS-coherence feature (the mean pairwise Pearson correlation of the physical RMS channels within the window) brings the total to 49. Feature names follow the pattern \texttt{\{family\}\_ch\{mean|max\}\_w\{mean|std|max|slope\}}, e.g.\ \texttt{rms\_chmax\_wslope}, plus the single \texttt{rms\_coherence}. A train-fit variance top-$k$ selection ($k \le 50$) then keeps $p < n$; every statistic is fit on the training split only, so no test-region information leaks into the features.

\begin{table}[t] \small
	\begin{center}
	\caption{The six per-snapshot base statistics from which the 49-dim invariant space is built. Each is computed per physical channel on the snapshot waveform $x$ (mean $\mu$, std $\sigma$), then channel-aggregated (mean, max) and window-summarised (mean, std, max, slope).}
	\label{tab:basestats}
	\begin{tabular}{l l l}
		\hline \hline
		\textbf{Statistic} & \textbf{Family} & \textbf{Formula} \\
		\hline
		RMS & rms & $\sqrt{\tfrac{1}{n}\sum_i x_i^2}$ \\
		Kurtosis & kurt & $\tfrac{1}{n}\sum_i (\tfrac{x_i-\mu}{\sigma})^4 - 3$ \\
		Skewness & skew & $\tfrac{1}{n}\sum_i (\tfrac{x_i-\mu}{\sigma})^3$ \\
		Peak-to-peak & p2p & $\max_i x_i - \min_i x_i$ \\
		Crest factor & crest & $\max_i |x_i| / \text{RMS}$ \\
		Variance & var & $\tfrac{1}{n}\sum_i (x_i-\mu)^2$ \\
		\hline \hline
	\end{tabular}
	\end{center}
\end{table}

\FloatBarrier
\section*{Appendix C: IMS Test-3 Relabeling Justification}
A subtlety we correct here is the failure-label of the third test: an earlier label placed failure on 2004-04-08, leaving 21.5\% of the recorded data after the nominal ``failure''---impossible for a run-to-failure series---so under that label the third test is a guaranteed 100\% miss for every detector. We relabel it to the last recorded sample (2004-04-18 02:42), at which the mean-RMS health indicator has ramped from its flat baseline (${\sim}0.068$) to ${\sim}0.228$ before the final machine-off reading---within the flat normal-region baseline band at the old label, but well above the $4\sigma$ onset band at the corrected one. The relabeling therefore moves the third test in the direction of harder, not easier, detection: a previously guaranteed 100\% miss becomes an evaluable run, and because the aggregate-vs-decimate contrast is computed independently within each run, the relabeling cannot artificially inflate it. We present this as an author relabeling based on the health-indicator trajectory, not an independently verified ground truth: we are not aware of a published benchmark that corrects this label, so we do not claim it as a community-established correction. It is nonetheless not a tuning choice: it only affects the single hardest run, and it makes that run pass rather than improving any aggregate-vs-decimate contrast (the contrast is computed within each run). It is the labeling we use throughout.

\FloatBarrier
\section*{Appendix D: Mechanism Analysis and ONGC Case Study}

\subsection*{D.1 Generic Smoothing vs.\ Historian Averaging}
The IMS trend admits a clean candidate mechanism: bin-averaging is a low-pass filter that suppresses single-window noise in the health statistic, letting the alarm threshold be crossed earlier and held more persistently than under decimation. If that is the mechanism, two predictions follow, and we test both directly on IMS (controlling window content at sampling factor 5 so only the smoothing differs).

\emph{Prediction 1: the aggregate advantage should grow with noise.} We add zero-mean Gaussian noise to the per-window health statistic at signal-to-noise ratios of 30, 20, and 10~dB and recompute the (aggregate $-$ decimate) lead time. To avoid the pseudoreplication we corrected in the main analysis, we collapse the four magnitude-monitoring SPC charts to one mean difference per run and treat the run as the unit of inference ($n=3$ per level). The mean advantage rises monotonically as noise increases (Table~\ref{tab:noise}): from $-0.35$~h at the native noise level to $+6.6$~h at 10~dB, and at the two highest noise levels all three runs share the positive sign. The lone large negative in the table, the $-7.94$~h native-noise entry, is IMS test 1---the abrupt near-coincident-onset run whose degradation begins only ${\sim}12$~h before failure, the same run that violates the conformal guarantee (Table~\ref{tab:conformal}) and is unwarnable in the spectral ablation (Table~\ref{tab:spectral}). This is the qualitative signature of a low-pass mechanism---the noisier the health signal, the more bin-averaging helps. We state its weight honestly: at $n=3$ this is a directional trend, not a significant effect (the exact sign test floors at $p=0.25$, and at 30~dB the per-run signs are mixed). It corroborates the proposed mechanism; it does not, and at this sample size cannot, establish it statistically.

\emph{Prediction 2: any smoother on the decimated stream should recover the advantage.} If the benefit is generic smoothing rather than the historian's averaging specifically, then applying an ordinary smoother to the decimated stream should match bin-averaging. We compare a median filter, a moving average, a constant-velocity Kalman filter, and a wavelet denoiser against raw aggregation and raw decimation (Table~\ref{tab:denoise}). The moving average and Kalman filter recover essentially the full aggregate lead time (76.0 and 75.7~h vs.\ aggregation's 75.7~h), while the median and wavelet filters fall slightly short. The effect is therefore generic smoothing of the health signal, not a property unique to historian bin-averaging---an operator who decimates and then lightly smooths recovers the same warning time. The linear smoothers (moving average, Kalman) recover the advantage but the wavelet denoiser does not, dropping to 6/12 valid alarms---below even raw decimation---because its soft-thresholding of detail coefficients perturbs the pre-onset baseline enough to lift the pre-onset false-alarm rate past the budget. The lesson sharpens rather than breaks: smoothing that preserves the normal baseline recovers the lead time, while smoothing that distorts it can cost validity. The one-alarm gap between raw decimation (8/12) and aggregation (7/12) at statistically indistinguishable lead times (76.0 vs.\ 75.7~h) is sampling noise at $n=12$, not evidence that decimation is preferable; it is exactly the non-destruction headline restated.

\begin{table}[t] \small
	\begin{center}
	\caption{IMS noise injection, run-level (factor 5): the four magnitude SPC charts collapsed to one mean aggregate $-$ decimate lead-time difference per run ($n=3$). Per-run values in order (test 1, test 2, test 3). The mean advantage grows monotonically as SNR falls---the signature of a low-pass mechanism---but at $n=3$ this is a directional trend that cannot be significance-tested.}
	\label{tab:noise}
	\begin{tabular}{l l r c}
		\hline \hline
		\textbf{Noise (SNR)} & \textbf{Per-run [t1,t2,t3]} & \textbf{Mean} & \textbf{Sign} \\
		\hline
		None (native) & $-7.94, +2.62, +4.28$ & $-0.35$ & 2/3 \\
		30 dB & $-0.44, -2.59, +5.95$ & $+0.97$ & 1/3 \\
		20 dB & $+10.60, +0.53, +0.12$ & $+3.75$ & 3/3 \\
		10 dB & $+14.99, +0.53, +4.28$ & $+6.60$ & 3/3 \\
		\hline \hline
	\end{tabular}
	\end{center}
\end{table}

\begin{table}[t] \small
	\begin{center}
	\caption{IMS denoiser comparison (factor 5, four magnitude charts collapsed to one mean per run, $n=3$): mean lead time when a smoother is applied to the decimated stream, against raw aggregation and raw decimation. A moving average or Kalman filter recovers the aggregate lead time, showing the benefit is generic smoothing.}
	\label{tab:denoise}
	\begin{tabular}{l r r}
		\hline \hline
		\textbf{Stream / denoiser} & \textbf{Mean lead (h)} & \textbf{Valid (/12)} \\
		\hline
		Aggregate (historian) & 75.7 & 7/12 \\
		Decimate (raw) & 76.0 & 8/12 \\
		Decimate + median & 69.6 & 8/12 \\
		Decimate + moving avg & 76.0 & 8/12 \\
		Decimate + Kalman & 75.7 & 8/12 \\
		Decimate + wavelet & 73.2 & 6/12 \\
		\hline \hline
	\end{tabular}
	\end{center}
\end{table}

\subsection*{D.2 Single-Asset ONGC Case Study}
The ONGC record is a real Solar-Turbine low-pressure-compressor (LPC) dataset: four vibration channels logged at a 10~s SCADA rate over ${\sim}5$ days (42,698 rows), ending in a genuine operator shutdown on 2023-11-13 09:46. This is authentic historian data with a real, documented event---the maintenance record confirms the shutdown followed a rolling-element bearing fault on the LPC---and the data-derived degradation onset (2023-11-13 03:44) precedes the operator's action by roughly six hours. ONGC is the only dataset that requires no simulated coarsening to be SCADA-rate: it is the historian.

On genuine 10~s industrial SCADA data ending in a real operator shutdown, every detector achieves a long warning: $3\sigma$, EWMA, Hotelling $T^2$, and Isolation Forest each alarm roughly 35~h before the labeled failure, and even the naive RMS-trend baseline gives ${\sim}22.8$~h. More to the point of this paper, the aggregate-vs-decimate difference is at most about a minute for every detector (Table~\ref{tab:ongc}). Because ONGC is a single asset ($n=1$), we treat it as a descriptive case study and draw no statistical inference from it. On data that is already a real historian stream, the choice of logging mechanism does not materially change warning time. This is the cleanest possible illustration of the practical question---it uses no simulated coarsening---and it is consistent with ``averaging is not the enemy.'' The conformal calibration on this stream (Figure~\ref{fig:conf_ongc}) tracks the diagonal closely, consistent with a mostly-stationary normal region.

\begin{table}[H] \small
	\begin{center}
	\caption{ONGC real-turbine SCADA (single-asset case study, $n=1$; no statistical inference): descriptive aggregate $-$ decimate lead-time difference, median over the five sampling factors. All differences are about a minute or less.}
	\label{tab:ongc}
	\begin{tabular}{l r}
		\hline \hline
		\textbf{Method} & \textbf{Median $\Delta$ (min)} \\
		\hline
		$3\sigma$ & $+0.5$ \\
		EWMA & $-1.1$ \\
		Hotelling $T^2$ & $+0.8$ \\
		Isolation Forest & $+1.2$ \\
		RMS-trend & 0.0 \\
		\hline \hline
	\end{tabular}
	\end{center}
\end{table}

\begin{figure}[H]
\centering
\includegraphics[width=\columnwidth]{fig_conf_ongc}
\caption{Conformal calibration on the ONGC turbine ($n=1$ case study): empirical pre-onset FAR vs target $\alpha$. The real historian stream tracks the diagonal closely, a modest 1.5--2$\times$ above target, consistent with a mostly-stationary normal region.}
\label{fig:conf_ongc}
\end{figure}

\end{document}
